% Insert with % Insert with % Insert with % Insert with \input{PSD_LIKELIHOOD_FINDINGS} after the Diagnostics and Interpretation Strategy section.
\section{Fixed-PSD Likelihood Audit and Interpretation}

\paragraph{Scope.}
The tables retain historical fixed-PSD audit results. They are not a current
released-map convergence certificate or evidence that BayesLine outperforms
Welch; later input-release qualifications below remain essential.

An additional audit was performed for GRB170817529/GW170817 and
GRB240715239 to make clear which likelihood quantities can be compared across
PSD choices. The sampling likelihood originally used in the pipeline is a
CBC-versus-noise likelihood ratio. For a fixed data segment and a fixed PSD,
this is the appropriate quantity for comparing the conditioned and
unconditioned sky priors: the omitted noise-only term is common to both
models and therefore cancels in $\Delta\log Z$. It is not, however, an
absolute likelihood that can be compared directly between Welch and BayesLine
PSDs.

For the audit, the fixed-PSD one-sided Whittle noise term was evaluated
post-hoc using the same $128\,\mathrm{s}$ Tukey-windowed data and frequency
grid as the sampler:
\[
\log p(d\mid N,S_n) = \sum_k \left[
\log\!\left(\frac{2\Delta f}{\pi S_n(f_k)}\right)
- \frac{2\Delta f\lvert\tilde d_k\rvert^2}{S_n(f_k)}
\right].
\]
This gives the bookkeeping identity
\[
\log Z_{\rm CBC} = \log p(d\mid N,S_n)
+ \log\mathcal{B}_{\rm CBC/N}.
\]
The first term depends on both the PSD determinant and the residual weighting,
so it is not a universal additive normalization. Changing the PSD changes the
relative weights of frequency bins and can consequently change the likelihood
geometry, posterior, and Bayes factor.

\begin{table}[ht]
\centering
\caption{Post-hoc fixed-PSD CBC-versus-noise results. Values are natural
logarithms; uncertainties are nested-sampling bootstrap estimates. The Welch
rows use the established $32\,\mathrm{s}$, $50\%$-overlap median estimator.}
\begin{tabular}{llrrr}
\toprule
\textbf{Event} & \textbf{PSD} & \textbf{$\log\mathcal{B}_{\rm CBC/N}^{\rm cond}$}
& \textbf{$\log\mathcal{B}_{\rm CBC/N}^{\rm uncond}$}
& \textbf{$\Delta\log Z$}\\
\midrule
GRB170817529 & BayesLine & $498.67 \pm 0.31$ & $491.80 \pm 0.34$ & $6.88 \pm 0.46$\\
GRB170817529 & Welch & $522.70 \pm 0.31$ & $517.48 \pm 0.33$ & $5.22 \pm 0.45$\\
GRB240715239 & BayesLine & $8.91 \pm 0.23$ & $5.74 \pm 0.23$ & $3.17 \pm 0.33$\\
GRB240715239 & Welch & $15.46 \pm 0.23$ & $13.54 \pm 0.27$ & $1.92 \pm 0.35$\\
\bottomrule
\end{tabular}
\end{table}

\begin{table}[ht]
\centering
\small
\caption{Post-hoc normalized fixed-PSD likelihood statistics. The normalized
CBC evidence is $\log Z_{\rm CBC}=\log p(d\mid N,S_n)+
\log\mathcal{B}_{\rm CBC/N}$, using the fixed-PSD Whittle convention. Values
are natural logarithms; uncertainties are nested-sampling bootstrap estimates.}
\begin{tabular}{llrrrr}
\toprule
\textbf{Event} & \textbf{PSD/prior} & \textbf{$\log p(d\mid N,S_n)$}
& \textbf{$\log Z_{\rm CBC}$} & \textbf{Mean $\log L$}
& \textbf{$D_{\rm KL}$ (nats)}\\
\midrule
GRB170817529 & BayesLine, conditioned & $76513731.438$ & $76514230.110$
& $76514266.810 \pm 0.086$ & $36.699 \pm 0.294$\\
GRB170817529 & BayesLine, unconditioned & $76513731.438$ & $76514223.233$
& $76514266.781 \pm 0.074$ & $43.547 \pm 0.341$\\
GRB170817529 & Welch, conditioned & $76534785.939$ & $76535308.640$
& $76535345.285 \pm 0.082$ & $36.646 \pm 0.294$\\
GRB170817529 & Welch, unconditioned & $76534785.939$ & $76535303.423$
& $76535345.359 \pm 0.079$ & $41.936 \pm 0.329$\\
\addlinespace
GRB240715239 & BayesLine, conditioned & $51677741.063$ & $51677749.975$
& $51677771.095 \pm 0.260$ & $21.120 \pm 0.268$\\
GRB240715239 & BayesLine, unconditioned & $51677741.063$ & $51677746.807$
& $51677766.637 \pm 0.136$ & $19.830 \pm 0.189$\\
GRB240715239 & Welch, conditioned & $51658062.380$ & $51658077.842$
& $51658098.695 \pm 0.112$ & $20.853 \pm 0.216$\\
GRB240715239 & Welch, unconditioned & $51658062.380$ & $51658075.921$
& $51658103.335 \pm 0.172$ & $27.413 \pm 0.250$\\
\bottomrule
\end{tabular}
\end{table}

These results show genuine sensitivity to the fixed PSD, rather than evidence
for one missing scalar normalization. In particular, GRB240715239 retains a
positive conditioned-versus-unconditioned comparison under both PSD choices,
but its magnitude changes. GW170817 remains decisively CBC-supported under
both choices and remains a useful control. The very large absolute values of
$\log p(d\mid N,S_n)$ and $\log Z_{\rm CBC}$ should not be interpreted as
an astrophysical ranking of PSD methods: they are dominated by the many-bin
noise likelihood and depend on the chosen fixed-PSD convention.

\subsection{Why the Welch--BayesLine Residual Difference Is Not ``False Physics''}

\textbf{Historical interpretation, qualified by the original-input audit
of 1 November 2026 below.} The PSD curves were subsequently shown to use
different published strain releases. The discussion of estimator and window
effects here must therefore not be treated as an identical-data comparison.

The large residual seen in the Welch-minus-BayesLine score is not generated by
Welch as an astrophysical feature. It is a consequence of the two fixed PSD
curves assigning very different noise power to a narrow instrumental feature,
most visibly in V1 near $1875\,\mathrm{Hz}$ for GW170817. Inspection of the
Welch PSD shows the broad, approximately symmetric skirts expected from
spectral leakage around a narrow feature. Across those skirts, the fixed Welch
and BayesLine PSDs differ by orders of magnitude in some bins. Bins with low
observed Fourier power are consequently assigned very different Whittle
weights, $1/S_n(f_k)$, and can make a large cumulative contribution to the
residual score. Expressed as a Welch-minus-BayesLine score, this produces a
large positive V1 contribution.

This does \emph{not} establish that Welch is physically correct or that
BayesLine is physically false. It identifies a calibration and modelling
issue: a posterior-median, off-source BayesLine PSD can under-represent the
power of a narrow line in the particular on-source realization, while a
median Welch estimate can absorb that line power more conservatively. A
held-out calibration test supports this interpretation. BayesLine is close to
the exponential-noise reference in the broad H1/L1/V1 distributions, while
the $32\,\mathrm{s}$ Welch estimate has a heavier upper tail, consistent with
localized excess power relative to its fixed PSD. Neither observation alone
selects a complete physical noise model for the on-source data.

Welch has no explicit transient-glitch model analogous to the wavelet
component available in BayesWave alongside its BayesLine PSD model. Median
averaging can reduce the influence of outlying periodograms, but it does not
identify, reconstruct, or subtract a glitch. This distinction must be kept
separate from spectral leakage caused by the finite-data window.

\subsection{Confirmation and Masking Tests}

The origin of the discrepancy was tested in three complementary ways. First,
the detector-resolved normalized score identified V1 as the only detector
with a large positive Welch-minus-BayesLine contribution. Second, a
frequency-resolved sum showed that at least $88\%$ of that V1 discrepancy is
concentrated in the narrow interval $1871$--$1879\,\mathrm{Hz}$, rather than
being distributed across the astrophysically relevant inspiral band. Third,
BayesWave/BayesLine was rerun with its glitch-wavelet component disabled. The no-glitch PSD did
not show a material drift from the standard glitch-modelled BayesLine PSD in
this region. The recorded comparison therefore showed no indication that
allocating transient power to glitch wavelets was responsible for the
discrepancy with Welch. It supports a narrow-feature/PSD-estimation or
window-response explanation rather than a difference in glitch treatment;
it does not, by itself, prove that no transient was present. This V1
high-frequency control is distinct from the known on-source L1 glitch in
GW170817, which an off-source PSD analysis does not remove.

\begin{table}[ht]
\centering
\small
\caption{GW170817 detector-resolved normalized fixed-PSD score difference,
$\ell^{\rm Welch}-\ell^{\rm BL}$, before any masking. These are absolute
fixed-PSD likelihood scores, not CBC Bayes factors.}
\begin{tabular}{lr}
\toprule
\textbf{Detector} & \textbf{$\ell^{\rm Welch}-\ell^{\rm BL}$ (nats)}\\
\midrule
H1 & $-1941.71$\\
L1 & $-4770.43$\\
V1 & $+27766.64$\\
\midrule
Network & $+21054.50$\\
\bottomrule
\end{tabular}
\end{table}

The masking tests are decisive as a diagnostic of the source of the absolute
score difference:
\begin{itemize}
\item Masking only $1871$--$1879\,\mathrm{Hz}$ removes at least $88\%$ of the
V1 positive contribution. This localizes the effect to the narrow spectral
feature and its leakage structure.
\item Excluding V1 entirely changes the network Welch-minus-BayesLine
normalized log likelihood from $+21054.50$ to $-6712.14$ nats. The sign flip
shows that V1, rather than H1 or L1, controls the original positive network
difference.
\end{itemize}

Neither operation is a replacement for a detector-consistent Bayesian
analysis. Frequency masking changes the likelihood and must be applied to
both the CBC and noise hypotheses before recomputing an evidence. Similarly,
the V1-excluded number is an H1/L1 noise-score diagnostic; an H1/L1 CBC Bayes
factor or evidence requires rerunning the likelihood and nested sampler with
V1 absent.

\subsection{Window-Response and Line-Evolution Test}

A follow-up on 30 October 2026 tested the finite-window response separately
from evolution of the dominant V1 line. The $1024\,\mathrm{s}$ off-source
context was divided into eight non-overlapping $128\,\mathrm{s}$ blocks.
Only the first four blocks were used to train the predictions; the remaining
four blocks and the on-source interval were withheld. An independent fit
in each block was retained as a diagnostic control, not as a held-out
prediction. These are narrow-band tests, not new CBC evidence calculations.

The dominant line was fitted near $1875.39202\,\mathrm{Hz}$, with an
all-block fitted-frequency range of only $3.8\times10^{-5}\,\mathrm{Hz}$.
Its peak strain amplitude nevertheless decreased from approximately
$9.46\times10^{-21}$ in the earliest block to $4.28\times10^{-21}$
on-source, corresponding to about an $80\%$ decline in line power.
A log-linear amplitude trend fitted exclusively to the four training blocks
had an amplitude e-folding time of approximately $1389\,\mathrm{s}$.
It predicted the later fitted line powers within about $6\%$, and the
on-source fitted line power within about $9\%$. This demonstrates amplitude
evolution in the cached strain; it does not identify the physical cause or
establish a transient-glitch classification.

The forward prediction used a sinusoid with that training-only frequency
and amplitude trend, including its within-block envelope, plus an estimated
local continuum. Its exact phase-averaged response was computed through
the symmetric Tukey window used by the on-source FFT. This did not convolve
an already-windowed PSD a second time. Synthetic-tone recovery,
window-weighted power conservation, and agreement with the pipeline's
on-source FFT periodogram were checked.

For the actual on-source Tukey window, the observed-to-predicted integrated
power ratios were $0.913$ in the core and $1.044$ in the skirts. The core
was defined within $0.25\,\mathrm{Hz}$ of the training frequency; the skirts
were $0.25$--$4\,\mathrm{Hz}$ away. Across the four held-out off-source
blocks these ratios ranged from $0.938$ to $1.004$ in the core and from
$1.022$ to $1.142$ in the skirts. In contrast, the stationary-amplitude
prediction gave an on-source core ratio of only $0.269$: matching the window
alone did not account for the change in line strength.

Within the local cache, this establishes a finite-window response and
nonstationary line amplitude. It does not by itself establish their share
of the cross-release Welch--BayesLine discrepancy. Agreement in
integrated power is not a complete calibration: the on-source mean
per-bin skirt-power ratio was still about $1.69$, versus a median of
$0.95$, indicating remaining localized excess. A Hann-window diagnostic
exposed additional narrow structures, leaving skirt power about $14$ times
the single-line-plus-continuum prediction. Parameter uncertainty,
additional lines, and frequency-bin covariance remain unmodelled. Local
core-conditioned fits are explanatory controls, and the reported diagonal
noise-score shifts are not model-selection Bayes factors. The BayesLine PSD
was displayed only as an on-source reference, not as an independently
trained held-out model. Production windows, PSD settings, and queued jobs
were unchanged.

\subsection{Matched No-Glitch BayesLine Control}

A diagnostic on 30 October 2026 matched the nominal Welch settings to the
archived \texttt{noiseOnly}, \texttt{noClean}, \texttt{noWaveletPrior}
BayesLine run. Its analysis interval was GPS
$[1187008610.4,1187008738.4]$, with $4096\,\mathrm{Hz}$ sampling,
a $128\,\mathrm{s}$ FFT, $\Delta f=0.0078125\,\mathrm{Hz}$, and the
$23$--$2048\,\mathrm{Hz}$ band. A $0.4\,\mathrm{s}$ edge taper,
corresponding to Tukey $\alpha=0.00625$, reproduces the logged window
correction to its printed precision. The separate configuration option
\texttt{window=1.0} was not interpreted as a Tukey alpha.

One complete segment produces a single periodogram, not a variance-reduced
Welch average. Its normalization was checked against a direct FFT and
window-weighted Parseval identity. Raw posterior PSDs were converted by
$2/128$ and checked against the archived physical-units export; the same
conversion was applied to the archived starting PSD. BayesWavePost omits
Nyquist and rounds frequency labels, so ordered rows were checked against
the exact starting-PSD grid instead of interpolating the rounded labels.

The nominally matched local strain still disagreed strongly with the
BayesLine PSD. However, the archived no-glitch Fourier residual export
revealed a more important distinction: it was not numerically equivalent
to the local strain. Integrated data-to-BayesLine power ratios were
\[
\begin{array}{lrr}
 & \text{Local cache} & \text{Archived Fourier export}\\
\text{Core: within }0.25\,\mathrm{Hz} & 1312.61 & 0.693\\
\text{Skirts: }0.25\text{--}4\,\mathrm{Hz} & 10.993 & 0.900\\
\text{Background: }5\text{--}8\,\mathrm{Hz} & 1.036 & 0.971
\end{array}
\]
around the fitted $1875.3920\,\mathrm{Hz}$ line. The archived export had
about $1900$ times less integrated core power than the local cache.
Symmetric versus periodic Tukey windows changed local core power by only
about $3.1\times10^{-7}$ fractionally. The nearest local sample starts
$0.4004$ samples earlier than the requested GPS time; adjacent-sample
sensitivity controls were also recorded.

The executed run flags disable glitch, signal, and CBC components, making
its Fourier residual a useful processed-data control. Nevertheless, it is
an exported residual product, not an independently verified copy of the
unconditioned frame strain. These results therefore locate an unresolved
data-provenance, conditioning, or export-consistency issue; they do not
identify its cause, demonstrate deliberate line removal, or establish that
BayesLine intrinsically misses the line. The archive-only core and skirt
ratios are of order one, but this is an in-sample diagnostic, not noise-model
calibration. The earlier conclusion that matching nominal settings alone
tests estimator differences must be qualified by this processed-data check.

The full $1024\,\mathrm{s}$ initialization context was not reproduced:
its first $2.075\,\mathrm{s}$ are absent from the local cache. No new
BayesWave or CBC inference jobs were launched, and production settings were
unchanged. Numerical checks, source hashes, arrays, and plots are saved by
the \texttt{--compare-no-glitch} diagnostic in its separate
\texttt{matched\_no\_glitch} output directory.

\subsection{Original-Input Release Identification and L1 Glitch Status}

On 1 November 2026 the actual frame input was checked rather than relying
on the post-processing export or nominal estimator settings. The archived
H1/L1/V1 frame-cache records name the $4096\,\mathrm{Hz}$ GWTC-1
GW170817 event-v3 files beginning at GPS $1187006835$, with duration
$4096\,\mathrm{s}$. SHA256 hashes of all three local frames matched their
current official GWOSC download bytes exactly. Duplicate local frame
copies also matched. This authenticates the available original input;
it is not a recovery of the historical download log.

The local sampler cache was separately compared with the official
O1\_O2-Preliminary GW170817 event-v2 cleaned HDF5 products beginning at
GPS $1187007040$. \emph{Every one of the $4{,}784{,}128$ aligned samples
per detector matched exactly}, over the complete cached
$1168\,\mathrm{s}$ interval
$[1187007716.4746094,1187008884.4746094)$. All H1/L1/V1 cache copies
checked against the local historical archive also had identical file
hashes. Thus the cached strain is identified as the published v2 release,
not as an unspecified O2 bulk product. The collection code's
\texttt{fetch\_open\_data(version=2)} is consistent with this finding,
but the identification rests on sample equality, not the version argument.

For the nominal $128\,\mathrm{s}$ no-glitch control, the first frame
sample at or after the requested start is GPS $1187008610.4001465$.
Using the same aligned local samples, symmetric Tukey taper and
periodogram normalization gives original-frame/local-cache integrated
V1 power ratios
\[
\begin{array}{lr}
\text{Core: within }0.25\,\mathrm{Hz}\text{ of }1875.3920\,\mathrm{Hz}
 & 0.00053034\\
\text{Skirts: }0.25\text{--}4\,\mathrm{Hz} & 0.0822164\\
\text{Broadband: }100\text{--}1500\,\mathrm{Hz} & 0.936858.
\end{array}
\]
The original frame already has approximately $1886$ times less core power.
The archived V1 Fourier export closely reproduces that frame: the relative
complex residual is approximately $0.00119$, with fitted complex scale
$0.99999908+0.00000081i$. Its integrated power divided by the physical
frame periodogram is approximately $0.996093$, consistent with the
window-energy factor $\langle w^2\rangle$ omitted from the export's
convention. A one-sample difference between nearest-sample and
first-at-or-after-start extraction materially affects the complex phases,
but not the orders-of-magnitude power distinction.

\textbf{The large V1 discrepancy is therefore present in the input strain
before BayesLine fitting.} It is not evidence that BayesLine intrinsically
misses the strong line seen in a different release. The successful complex
match also rules out a wholly independent simulated-noise realization as
an adequate explanation for this export, despite the simulated-channel
labels in the BayesWavePost log. The exact remaining conditioning,
calibration and downsampling differences have not been isolated, and the
precise historical post-processing overwrite sequence is not inferred.

The known strong GW170817 glitch is a separate L1 issue. GWOSC documents
its time as GPS $1187008881.389$, about $1.1\,\mathrm{s}$ before merger,
and a brief digital-to-analog converter saturation of less than
$5\,\mathrm{ms}$. GWTC-1 event-v3 strain retains the glitch; the earlier
v2 release is explicitly documented as glitch-subtracted. Since the full
cached L1 samples exactly match v2, \emph{the local sampler cache already
uses the published glitch-subtracted L1 data}. A common diagnostic
$30$--$1000\,\mathrm{Hz}$ zero-phase bandpass shows the transient in the
original v3 frame and its suppression in the cache. This plot is a display
control, not a reconstruction-validation test or a new likelihood.
Unfiltered raw maxima alone are misleading because the v2 H1/L1 products
have documented excess noise below $20\,\mathrm{Hz}$. The v2 H1
last-$150\,\mathrm{s}$ cleaning-window warning concerns approximately
GPS $1187008938$ onward, outside the cached interval checked here.

The off-source BayesWave control ends at GPS $1187008738.4$, before the
L1 glitch. Off-source PSD cleaning neither explains the cross-detector
V1 difference nor establishes on-source glitch subtraction. In this case,
the cache's subtraction status is established independently by its release
content. Special glitch-subtracted products described by GWOSC also differ
in calibration and glitch model between publications; they must not be
interchanged merely because all are called cleaned data.

There is an additional frequency-band caveat. GWOSC's O2 technical notes
recommend the $16384\,\mathrm{Hz}$ products for studies around
$1700\,\mathrm{Hz}$ and above because the $4096\,\mathrm{Hz}$
anti-alias response corrupts the Nyquist region. The $1875\,\mathrm{Hz}$
diagnostic cannot be promoted to physical high-frequency calibration
without explicit release and filter checks. No frequency cut, data
replacement, gating, PSD regeneration or new inference job was applied.
The frozen calibration bundle remains a legacy mixed-release fixed-PSD
comparison; its historical unresolved-glitch caveat is now qualified by
this independent sample audit, not rewritten in its frozen configuration.

Earlier tables and score snapshots are preserved. The historical
``at least $88\%$'' localization claim above was not reproduced by a
later direct net-score integration, which gave approximately $67\%$
for the $1871$--$1879\,\mathrm{Hz}$ share of the net V1 shift; these
definitions or snapshots must not be combined. Neither those score tests
nor the present release-identification checks validate a detector-consistent
PSD likelihood or a completed GW170817 recovery. A fair estimator
comparison requires PSD estimation and on-source strain from a matched
release with explicit conditioning and a defensible analysis band.

The independent original-input audit saves official API snapshots,
file and sample hashes, aligned spectra, ratios and raw/bandpassed glitch
plots in a fresh timestamped directory. All $16$ hashed source/config/input
files in the frozen calibration bundle were checked without modification.
Sources are the GWOSC GWTC-1 documentation, the GW170817 event-v2/v3
release records, the dedicated GW170817 strain documentation and the O2
technical details; URLs and downloaded-file hashes are recorded in the
audit report. No legacy inference results were replaced.

\subsection{Top-Candidate Cross-Checks}

The same detector- and frequency-resolved normalized-score diagnostic was
applied to the two leading subthreshold candidates, GRB190728271 and
GRB240715239. Unlike GW170817, their discrepancies do not support a common
V1 feature near $1875\,\mathrm{Hz}$. Table~\ref{tab:top-candidate-psd-scores}
uses the sign convention $\ell^{\rm BL}-\ell^{\rm Welch}$; a positive value
therefore indicates a higher normalized fixed-PSD score for BayesLine.

\begin{table}[ht]
\centering
\small
\caption{Detector-resolved normalized fixed-PSD score differences for the
top candidates. The CBC entries use the conditioned maximum-log-likelihood
posterior waveform with the common coalescence phase profiled separately for
each PSD. These are predictive-score diagnostics, not Bayes factors.}
\label{tab:top-candidate-psd-scores}
\begin{tabular}{llrr}
\toprule
\textbf{Event} & \textbf{Detector} & \textbf{Noise-only}
& \textbf{CBC residual}\\
& & \textbf{$\ell^{\rm BL}-\ell^{\rm Welch}$} &
\textbf{$\ell^{\rm BL}-\ell^{\rm Welch}$}\\
\midrule
GRB190728271 & H1 & $+3031.38$ & $+3032.01$\\
GRB190728271 & L1 & $+4583.71$ & $+4581.30$\\
GRB190728271 & Network & $+7615.09$ & $+7613.31$\\
\addlinespace
GRB240715239 & L1 & $+4379.88$ & $+4377.77$\\
GRB240715239 & V1 & $+7616.29$ & $+7622.69$\\
GRB240715239 & Network & $+11996.16$ & $+12000.46$\\
\bottomrule
\end{tabular}
\end{table}

GRB190728271 was observed only in H1 and L1, so a V1 masking test is not
defined for that event. Its score difference is distributed across both
detectors. The largest $1\,\mathrm{Hz}$ contributions occur around
$1461$--$1465\,\mathrm{Hz}$ in H1 and $1022$--$1027\,\mathrm{Hz}$ in L1,
showing that its PSD sensitivity cannot be attributed to the GW170817 V1
feature.

For GRB240715239, dropping V1 leaves the L1-only normalized difference of
$+4379.88$ nats, rather than reversing the sign as it did for GW170817. V1 is
important, contributing $+7616.29$ nats, but L1 also contributes materially.
The targeted V1 interval $1871$--$1879\,\mathrm{Hz}$ contributes only
$+232.12$ nats, or $3.05\%$ of the V1 noise-only difference. Its largest V1
score differences instead lie predominantly between $31$ and
$64\,\mathrm{Hz}$. Thus, although both candidates are sensitive to the fixed
PSD choice, neither has evidence for the specific GW170817-like
$1875\,\mathrm{Hz}$ V1 discrepancy.

As above, detector removal or frequency masking changes the likelihood.
These results localize the fixed-PSD sensitivity, but they do not yield
masked CBC evidences or conditioned-versus-unconditioned Bayes factors.
Those quantities require detector- and mask-consistent likelihood evaluations
and nested-sampling reruns.

The conclusion is therefore methodological. Fixed-PSD evidence differences
should be treated as PSD-sensitivity diagnostics. They motivate inspections
of the affected detector and frequency range, repeat runs using posterior PSD
draws, off-source control studies, and software injections. A definitive
comparison would propagate PSD uncertainty and model non-stationary lines,
window-induced frequency correlations, and transient glitches jointly with the
CBC signal. Finally, removing V1 after the fact permits only an H1/L1
noise-score diagnostic; it does not produce an H1/L1 CBC Bayes factor or
evidence without rerunning the likelihood and nested sampler for that detector
network.

\subsection{Detector- and Frequency-Resolved Likelihood Digest}

For a detector $I$ and frequency bin $k$, define the normalized fixed-PSD
Whittle contribution for a waveform $h$ as
\[
\ell_{I,k}(h;S_n) =
\log\!\left(\frac{2\Delta f}{\pi S_{n,I}(f_k)}\right)
- \frac{2\Delta f\lvert\tilde d_{I,k}-\tilde h_{I,k}\rvert^2}
{S_{n,I}(f_k)}.
\]
The detector score and network score are respectively
\[
\ell_I(h;S_n) = \sum_{k\in[23,2048]~\mathrm{Hz}}\ell_{I,k}(h;S_n),
\qquad
\ell_{\rm net}(h;S_n) = \sum_I\ell_I(h;S_n).
\]
For the noise-only diagnostic, set $h=0$. For the CBC diagnostic, $h$ is a
specified posterior waveform, with a common coalescence phase profiled
separately for each fixed PSD. These are predictive scores for a selected
waveform; they are not nested-sampling evidences and should not replace the
evidence table above.

For direct comparison of the two PSD curves, the detector- and bin-resolved
differences are
\[
\delta\ell_I(h) = \ell_I(h;S_n^{\rm BL})-
\ell_I(h;S_n^{\rm Welch}),
\qquad
\delta\ell_{I,k}(h) =
\ell_{I,k}(h;S_n^{\rm BL})-
\ell_{I,k}(h;S_n^{\rm Welch}),
\]
with $\delta\ell_I=\sum_k\delta\ell_{I,k}$. A positive value means that
BayesLine assigns a higher normalized score; a negative value means that
Welch assigns a higher normalized score. This convention makes the V1
discrepancy transparent: the narrow bin near $1875\,\mathrm{Hz}$ gives a
large negative $\delta\ell_{\mathrm{V1},k}$ when the BayesLine PSD is low
relative to the observed on-source line power.

\begin{table}[ht]
\centering
\small
\caption{Detector-resolved format for the normalized Whittle audit. Each row
is obtained by summing the bin contributions in the stated detector over the
analysis band. The network row is the sum of detector rows. Report the values
from the audit output for the selected event and posterior waveform.}
\begin{tabular}{llrrrr}
\toprule
\textbf{Detector} & \textbf{Waveform} & \textbf{$\ell_I^{\rm Welch}$}
& \textbf{$\ell_I^{\rm BL}$} & \textbf{$\delta\ell_I$}
& \textbf{Interpretation}\\
\midrule
H1 & $h=0$ & \textit{audit value} & \textit{audit value}
& \textit{BL $-$ Welch} & noise-only score\\
L1 & $h=0$ & \textit{audit value} & \textit{audit value}
& \textit{BL $-$ Welch} & noise-only score\\
V1 & $h=0$ & \textit{audit value} & \textit{audit value}
& \textit{BL $-$ Welch} & inspect narrow lines\\
Network & $h=0$ & \textit{sum} & \textit{sum}
& \textit{sum} & detector sum\\
\addlinespace
H1 & $h=h_{\rm prof}$ & \textit{audit value} & \textit{audit value}
& \textit{BL $-$ Welch} & CBC residual score\\
L1 & $h=h_{\rm prof}$ & \textit{audit value} & \textit{audit value}
& \textit{BL $-$ Welch} & CBC residual score\\
V1 & $h=h_{\rm prof}$ & \textit{audit value} & \textit{audit value}
& \textit{BL $-$ Welch} & CBC residual score\\
Network & $h=h_{\rm prof}$ & \textit{sum} & \textit{sum}
& \textit{sum} & detector sum\\
\bottomrule
\end{tabular}
\end{table}

Rather than tabulating every $1/128\,\mathrm{Hz}$ Fourier bin, which would
obscure the result, the per-bin audit lists the bins or narrow bands with the
largest absolute score differences. The confirmed V1 feature can be reported
compactly as follows:
\begin{table}[ht]
\centering
\small
\caption{Largest frequency-resolved contributions to the fixed-PSD score
difference. Values are computed from $\delta\ell_{I,k}$; neighbouring bins
may be grouped only when their summed contribution is reported.}
\begin{tabular}{llrrrl}
\toprule
\textbf{Event} & \textbf{Detector} & \textbf{Frequency/band [Hz]}
& \textbf{$\delta\ell_{I,k}(0)$} & \textbf{$\delta\ell_{I,k}(h_{\rm prof})$}
& \textbf{Feature}\\
\midrule
GRB170817529 & V1 & $1871$--$1879$ &
$\sum_{k\in\mathrm{band}}\delta\ell_{\rm V1,k}(0)
\leq-0.88\lvert\delta\ell_{\rm V1}(0)\rvert$
& \textit{audit value} & narrow feature and leakage skirts\\
\textit{event} & \textit{IFO} & \textit{largest remaining bin/band}
& \textit{audit value} & \textit{audit value} & \textit{diagnostic label}\\
\textit{event} & \textit{IFO} & \textit{largest remaining bin/band}
& \textit{audit value} & \textit{audit value} & \textit{diagnostic label}\\
\bottomrule
\end{tabular}
\end{table}

The two tables should always be read together. A large detector total can be
broad-band, or it can be dominated by one narrow instrumental feature. The
per-bin table distinguishes these cases and prevents a narrow V1 line from
being misdescribed as broad astrophysical residual power. after the Diagnostics and Interpretation Strategy section.
\section{Fixed-PSD Likelihood Audit and Interpretation}

\paragraph{Scope.}
The tables retain historical fixed-PSD audit results. They are not a current
released-map convergence certificate or evidence that BayesLine outperforms
Welch; later input-release qualifications below remain essential.

An additional audit was performed for GRB170817529/GW170817 and
GRB240715239 to make clear which likelihood quantities can be compared across
PSD choices. The sampling likelihood originally used in the pipeline is a
CBC-versus-noise likelihood ratio. For a fixed data segment and a fixed PSD,
this is the appropriate quantity for comparing the conditioned and
unconditioned sky priors: the omitted noise-only term is common to both
models and therefore cancels in $\Delta\log Z$. It is not, however, an
absolute likelihood that can be compared directly between Welch and BayesLine
PSDs.

For the audit, the fixed-PSD one-sided Whittle noise term was evaluated
post-hoc using the same $128\,\mathrm{s}$ Tukey-windowed data and frequency
grid as the sampler:
\[
\log p(d\mid N,S_n) = \sum_k \left[
\log\!\left(\frac{2\Delta f}{\pi S_n(f_k)}\right)
- \frac{2\Delta f\lvert\tilde d_k\rvert^2}{S_n(f_k)}
\right].
\]
This gives the bookkeeping identity
\[
\log Z_{\rm CBC} = \log p(d\mid N,S_n)
+ \log\mathcal{B}_{\rm CBC/N}.
\]
The first term depends on both the PSD determinant and the residual weighting,
so it is not a universal additive normalization. Changing the PSD changes the
relative weights of frequency bins and can consequently change the likelihood
geometry, posterior, and Bayes factor.

\begin{table}[ht]
\centering
\caption{Post-hoc fixed-PSD CBC-versus-noise results. Values are natural
logarithms; uncertainties are nested-sampling bootstrap estimates. The Welch
rows use the established $32\,\mathrm{s}$, $50\%$-overlap median estimator.}
\begin{tabular}{llrrr}
\toprule
\textbf{Event} & \textbf{PSD} & \textbf{$\log\mathcal{B}_{\rm CBC/N}^{\rm cond}$}
& \textbf{$\log\mathcal{B}_{\rm CBC/N}^{\rm uncond}$}
& \textbf{$\Delta\log Z$}\\
\midrule
GRB170817529 & BayesLine & $498.67 \pm 0.31$ & $491.80 \pm 0.34$ & $6.88 \pm 0.46$\\
GRB170817529 & Welch & $522.70 \pm 0.31$ & $517.48 \pm 0.33$ & $5.22 \pm 0.45$\\
GRB240715239 & BayesLine & $8.91 \pm 0.23$ & $5.74 \pm 0.23$ & $3.17 \pm 0.33$\\
GRB240715239 & Welch & $15.46 \pm 0.23$ & $13.54 \pm 0.27$ & $1.92 \pm 0.35$\\
\bottomrule
\end{tabular}
\end{table}

\begin{table}[ht]
\centering
\small
\caption{Post-hoc normalized fixed-PSD likelihood statistics. The normalized
CBC evidence is $\log Z_{\rm CBC}=\log p(d\mid N,S_n)+
\log\mathcal{B}_{\rm CBC/N}$, using the fixed-PSD Whittle convention. Values
are natural logarithms; uncertainties are nested-sampling bootstrap estimates.}
\begin{tabular}{llrrrr}
\toprule
\textbf{Event} & \textbf{PSD/prior} & \textbf{$\log p(d\mid N,S_n)$}
& \textbf{$\log Z_{\rm CBC}$} & \textbf{Mean $\log L$}
& \textbf{$D_{\rm KL}$ (nats)}\\
\midrule
GRB170817529 & BayesLine, conditioned & $76513731.438$ & $76514230.110$
& $76514266.810 \pm 0.086$ & $36.699 \pm 0.294$\\
GRB170817529 & BayesLine, unconditioned & $76513731.438$ & $76514223.233$
& $76514266.781 \pm 0.074$ & $43.547 \pm 0.341$\\
GRB170817529 & Welch, conditioned & $76534785.939$ & $76535308.640$
& $76535345.285 \pm 0.082$ & $36.646 \pm 0.294$\\
GRB170817529 & Welch, unconditioned & $76534785.939$ & $76535303.423$
& $76535345.359 \pm 0.079$ & $41.936 \pm 0.329$\\
\addlinespace
GRB240715239 & BayesLine, conditioned & $51677741.063$ & $51677749.975$
& $51677771.095 \pm 0.260$ & $21.120 \pm 0.268$\\
GRB240715239 & BayesLine, unconditioned & $51677741.063$ & $51677746.807$
& $51677766.637 \pm 0.136$ & $19.830 \pm 0.189$\\
GRB240715239 & Welch, conditioned & $51658062.380$ & $51658077.842$
& $51658098.695 \pm 0.112$ & $20.853 \pm 0.216$\\
GRB240715239 & Welch, unconditioned & $51658062.380$ & $51658075.921$
& $51658103.335 \pm 0.172$ & $27.413 \pm 0.250$\\
\bottomrule
\end{tabular}
\end{table}

These results show genuine sensitivity to the fixed PSD, rather than evidence
for one missing scalar normalization. In particular, GRB240715239 retains a
positive conditioned-versus-unconditioned comparison under both PSD choices,
but its magnitude changes. GW170817 remains decisively CBC-supported under
both choices and remains a useful control. The very large absolute values of
$\log p(d\mid N,S_n)$ and $\log Z_{\rm CBC}$ should not be interpreted as
an astrophysical ranking of PSD methods: they are dominated by the many-bin
noise likelihood and depend on the chosen fixed-PSD convention.

\subsection{Why the Welch--BayesLine Residual Difference Is Not ``False Physics''}

\textbf{Historical interpretation, qualified by the original-input audit
of 1 November 2026 below.} The PSD curves were subsequently shown to use
different published strain releases. The discussion of estimator and window
effects here must therefore not be treated as an identical-data comparison.

The large residual seen in the Welch-minus-BayesLine score is not generated by
Welch as an astrophysical feature. It is a consequence of the two fixed PSD
curves assigning very different noise power to a narrow instrumental feature,
most visibly in V1 near $1875\,\mathrm{Hz}$ for GW170817. Inspection of the
Welch PSD shows the broad, approximately symmetric skirts expected from
spectral leakage around a narrow feature. Across those skirts, the fixed Welch
and BayesLine PSDs differ by orders of magnitude in some bins. Bins with low
observed Fourier power are consequently assigned very different Whittle
weights, $1/S_n(f_k)$, and can make a large cumulative contribution to the
residual score. Expressed as a Welch-minus-BayesLine score, this produces a
large positive V1 contribution.

This does \emph{not} establish that Welch is physically correct or that
BayesLine is physically false. It identifies a calibration and modelling
issue: a posterior-median, off-source BayesLine PSD can under-represent the
power of a narrow line in the particular on-source realization, while a
median Welch estimate can absorb that line power more conservatively. A
held-out calibration test supports this interpretation. BayesLine is close to
the exponential-noise reference in the broad H1/L1/V1 distributions, while
the $32\,\mathrm{s}$ Welch estimate has a heavier upper tail, consistent with
localized excess power relative to its fixed PSD. Neither observation alone
selects a complete physical noise model for the on-source data.

Welch has no explicit transient-glitch model analogous to the wavelet
component available in BayesWave alongside its BayesLine PSD model. Median
averaging can reduce the influence of outlying periodograms, but it does not
identify, reconstruct, or subtract a glitch. This distinction must be kept
separate from spectral leakage caused by the finite-data window.

\subsection{Confirmation and Masking Tests}

The origin of the discrepancy was tested in three complementary ways. First,
the detector-resolved normalized score identified V1 as the only detector
with a large positive Welch-minus-BayesLine contribution. Second, a
frequency-resolved sum showed that at least $88\%$ of that V1 discrepancy is
concentrated in the narrow interval $1871$--$1879\,\mathrm{Hz}$, rather than
being distributed across the astrophysically relevant inspiral band. Third,
BayesWave/BayesLine was rerun with its glitch-wavelet component disabled. The no-glitch PSD did
not show a material drift from the standard glitch-modelled BayesLine PSD in
this region. The recorded comparison therefore showed no indication that
allocating transient power to glitch wavelets was responsible for the
discrepancy with Welch. It supports a narrow-feature/PSD-estimation or
window-response explanation rather than a difference in glitch treatment;
it does not, by itself, prove that no transient was present. This V1
high-frequency control is distinct from the known on-source L1 glitch in
GW170817, which an off-source PSD analysis does not remove.

\begin{table}[ht]
\centering
\small
\caption{GW170817 detector-resolved normalized fixed-PSD score difference,
$\ell^{\rm Welch}-\ell^{\rm BL}$, before any masking. These are absolute
fixed-PSD likelihood scores, not CBC Bayes factors.}
\begin{tabular}{lr}
\toprule
\textbf{Detector} & \textbf{$\ell^{\rm Welch}-\ell^{\rm BL}$ (nats)}\\
\midrule
H1 & $-1941.71$\\
L1 & $-4770.43$\\
V1 & $+27766.64$\\
\midrule
Network & $+21054.50$\\
\bottomrule
\end{tabular}
\end{table}

The masking tests are decisive as a diagnostic of the source of the absolute
score difference:
\begin{itemize}
\item Masking only $1871$--$1879\,\mathrm{Hz}$ removes at least $88\%$ of the
V1 positive contribution. This localizes the effect to the narrow spectral
feature and its leakage structure.
\item Excluding V1 entirely changes the network Welch-minus-BayesLine
normalized log likelihood from $+21054.50$ to $-6712.14$ nats. The sign flip
shows that V1, rather than H1 or L1, controls the original positive network
difference.
\end{itemize}

Neither operation is a replacement for a detector-consistent Bayesian
analysis. Frequency masking changes the likelihood and must be applied to
both the CBC and noise hypotheses before recomputing an evidence. Similarly,
the V1-excluded number is an H1/L1 noise-score diagnostic; an H1/L1 CBC Bayes
factor or evidence requires rerunning the likelihood and nested sampler with
V1 absent.

\subsection{Window-Response and Line-Evolution Test}

A follow-up on 30 October 2026 tested the finite-window response separately
from evolution of the dominant V1 line. The $1024\,\mathrm{s}$ off-source
context was divided into eight non-overlapping $128\,\mathrm{s}$ blocks.
Only the first four blocks were used to train the predictions; the remaining
four blocks and the on-source interval were withheld. An independent fit
in each block was retained as a diagnostic control, not as a held-out
prediction. These are narrow-band tests, not new CBC evidence calculations.

The dominant line was fitted near $1875.39202\,\mathrm{Hz}$, with an
all-block fitted-frequency range of only $3.8\times10^{-5}\,\mathrm{Hz}$.
Its peak strain amplitude nevertheless decreased from approximately
$9.46\times10^{-21}$ in the earliest block to $4.28\times10^{-21}$
on-source, corresponding to about an $80\%$ decline in line power.
A log-linear amplitude trend fitted exclusively to the four training blocks
had an amplitude e-folding time of approximately $1389\,\mathrm{s}$.
It predicted the later fitted line powers within about $6\%$, and the
on-source fitted line power within about $9\%$. This demonstrates amplitude
evolution in the cached strain; it does not identify the physical cause or
establish a transient-glitch classification.

The forward prediction used a sinusoid with that training-only frequency
and amplitude trend, including its within-block envelope, plus an estimated
local continuum. Its exact phase-averaged response was computed through
the symmetric Tukey window used by the on-source FFT. This did not convolve
an already-windowed PSD a second time. Synthetic-tone recovery,
window-weighted power conservation, and agreement with the pipeline's
on-source FFT periodogram were checked.

For the actual on-source Tukey window, the observed-to-predicted integrated
power ratios were $0.913$ in the core and $1.044$ in the skirts. The core
was defined within $0.25\,\mathrm{Hz}$ of the training frequency; the skirts
were $0.25$--$4\,\mathrm{Hz}$ away. Across the four held-out off-source
blocks these ratios ranged from $0.938$ to $1.004$ in the core and from
$1.022$ to $1.142$ in the skirts. In contrast, the stationary-amplitude
prediction gave an on-source core ratio of only $0.269$: matching the window
alone did not account for the change in line strength.

Within the local cache, this establishes a finite-window response and
nonstationary line amplitude. It does not by itself establish their share
of the cross-release Welch--BayesLine discrepancy. Agreement in
integrated power is not a complete calibration: the on-source mean
per-bin skirt-power ratio was still about $1.69$, versus a median of
$0.95$, indicating remaining localized excess. A Hann-window diagnostic
exposed additional narrow structures, leaving skirt power about $14$ times
the single-line-plus-continuum prediction. Parameter uncertainty,
additional lines, and frequency-bin covariance remain unmodelled. Local
core-conditioned fits are explanatory controls, and the reported diagonal
noise-score shifts are not model-selection Bayes factors. The BayesLine PSD
was displayed only as an on-source reference, not as an independently
trained held-out model. Production windows, PSD settings, and queued jobs
were unchanged.

\subsection{Matched No-Glitch BayesLine Control}

A diagnostic on 30 October 2026 matched the nominal Welch settings to the
archived \texttt{noiseOnly}, \texttt{noClean}, \texttt{noWaveletPrior}
BayesLine run. Its analysis interval was GPS
$[1187008610.4,1187008738.4]$, with $4096\,\mathrm{Hz}$ sampling,
a $128\,\mathrm{s}$ FFT, $\Delta f=0.0078125\,\mathrm{Hz}$, and the
$23$--$2048\,\mathrm{Hz}$ band. A $0.4\,\mathrm{s}$ edge taper,
corresponding to Tukey $\alpha=0.00625$, reproduces the logged window
correction to its printed precision. The separate configuration option
\texttt{window=1.0} was not interpreted as a Tukey alpha.

One complete segment produces a single periodogram, not a variance-reduced
Welch average. Its normalization was checked against a direct FFT and
window-weighted Parseval identity. Raw posterior PSDs were converted by
$2/128$ and checked against the archived physical-units export; the same
conversion was applied to the archived starting PSD. BayesWavePost omits
Nyquist and rounds frequency labels, so ordered rows were checked against
the exact starting-PSD grid instead of interpolating the rounded labels.

The nominally matched local strain still disagreed strongly with the
BayesLine PSD. However, the archived no-glitch Fourier residual export
revealed a more important distinction: it was not numerically equivalent
to the local strain. Integrated data-to-BayesLine power ratios were
\[
\begin{array}{lrr}
 & \text{Local cache} & \text{Archived Fourier export}\\
\text{Core: within }0.25\,\mathrm{Hz} & 1312.61 & 0.693\\
\text{Skirts: }0.25\text{--}4\,\mathrm{Hz} & 10.993 & 0.900\\
\text{Background: }5\text{--}8\,\mathrm{Hz} & 1.036 & 0.971
\end{array}
\]
around the fitted $1875.3920\,\mathrm{Hz}$ line. The archived export had
about $1900$ times less integrated core power than the local cache.
Symmetric versus periodic Tukey windows changed local core power by only
about $3.1\times10^{-7}$ fractionally. The nearest local sample starts
$0.4004$ samples earlier than the requested GPS time; adjacent-sample
sensitivity controls were also recorded.

The executed run flags disable glitch, signal, and CBC components, making
its Fourier residual a useful processed-data control. Nevertheless, it is
an exported residual product, not an independently verified copy of the
unconditioned frame strain. These results therefore locate an unresolved
data-provenance, conditioning, or export-consistency issue; they do not
identify its cause, demonstrate deliberate line removal, or establish that
BayesLine intrinsically misses the line. The archive-only core and skirt
ratios are of order one, but this is an in-sample diagnostic, not noise-model
calibration. The earlier conclusion that matching nominal settings alone
tests estimator differences must be qualified by this processed-data check.

The full $1024\,\mathrm{s}$ initialization context was not reproduced:
its first $2.075\,\mathrm{s}$ are absent from the local cache. No new
BayesWave or CBC inference jobs were launched, and production settings were
unchanged. Numerical checks, source hashes, arrays, and plots are saved by
the \texttt{--compare-no-glitch} diagnostic in its separate
\texttt{matched\_no\_glitch} output directory.

\subsection{Original-Input Release Identification and L1 Glitch Status}

On 1 November 2026 the actual frame input was checked rather than relying
on the post-processing export or nominal estimator settings. The archived
H1/L1/V1 frame-cache records name the $4096\,\mathrm{Hz}$ GWTC-1
GW170817 event-v3 files beginning at GPS $1187006835$, with duration
$4096\,\mathrm{s}$. SHA256 hashes of all three local frames matched their
current official GWOSC download bytes exactly. Duplicate local frame
copies also matched. This authenticates the available original input;
it is not a recovery of the historical download log.

The local sampler cache was separately compared with the official
O1\_O2-Preliminary GW170817 event-v2 cleaned HDF5 products beginning at
GPS $1187007040$. \emph{Every one of the $4{,}784{,}128$ aligned samples
per detector matched exactly}, over the complete cached
$1168\,\mathrm{s}$ interval
$[1187007716.4746094,1187008884.4746094)$. All H1/L1/V1 cache copies
checked against the local historical archive also had identical file
hashes. Thus the cached strain is identified as the published v2 release,
not as an unspecified O2 bulk product. The collection code's
\texttt{fetch\_open\_data(version=2)} is consistent with this finding,
but the identification rests on sample equality, not the version argument.

For the nominal $128\,\mathrm{s}$ no-glitch control, the first frame
sample at or after the requested start is GPS $1187008610.4001465$.
Using the same aligned local samples, symmetric Tukey taper and
periodogram normalization gives original-frame/local-cache integrated
V1 power ratios
\[
\begin{array}{lr}
\text{Core: within }0.25\,\mathrm{Hz}\text{ of }1875.3920\,\mathrm{Hz}
 & 0.00053034\\
\text{Skirts: }0.25\text{--}4\,\mathrm{Hz} & 0.0822164\\
\text{Broadband: }100\text{--}1500\,\mathrm{Hz} & 0.936858.
\end{array}
\]
The original frame already has approximately $1886$ times less core power.
The archived V1 Fourier export closely reproduces that frame: the relative
complex residual is approximately $0.00119$, with fitted complex scale
$0.99999908+0.00000081i$. Its integrated power divided by the physical
frame periodogram is approximately $0.996093$, consistent with the
window-energy factor $\langle w^2\rangle$ omitted from the export's
convention. A one-sample difference between nearest-sample and
first-at-or-after-start extraction materially affects the complex phases,
but not the orders-of-magnitude power distinction.

\textbf{The large V1 discrepancy is therefore present in the input strain
before BayesLine fitting.} It is not evidence that BayesLine intrinsically
misses the strong line seen in a different release. The successful complex
match also rules out a wholly independent simulated-noise realization as
an adequate explanation for this export, despite the simulated-channel
labels in the BayesWavePost log. The exact remaining conditioning,
calibration and downsampling differences have not been isolated, and the
precise historical post-processing overwrite sequence is not inferred.

The known strong GW170817 glitch is a separate L1 issue. GWOSC documents
its time as GPS $1187008881.389$, about $1.1\,\mathrm{s}$ before merger,
and a brief digital-to-analog converter saturation of less than
$5\,\mathrm{ms}$. GWTC-1 event-v3 strain retains the glitch; the earlier
v2 release is explicitly documented as glitch-subtracted. Since the full
cached L1 samples exactly match v2, \emph{the local sampler cache already
uses the published glitch-subtracted L1 data}. A common diagnostic
$30$--$1000\,\mathrm{Hz}$ zero-phase bandpass shows the transient in the
original v3 frame and its suppression in the cache. This plot is a display
control, not a reconstruction-validation test or a new likelihood.
Unfiltered raw maxima alone are misleading because the v2 H1/L1 products
have documented excess noise below $20\,\mathrm{Hz}$. The v2 H1
last-$150\,\mathrm{s}$ cleaning-window warning concerns approximately
GPS $1187008938$ onward, outside the cached interval checked here.

The off-source BayesWave control ends at GPS $1187008738.4$, before the
L1 glitch. Off-source PSD cleaning neither explains the cross-detector
V1 difference nor establishes on-source glitch subtraction. In this case,
the cache's subtraction status is established independently by its release
content. Special glitch-subtracted products described by GWOSC also differ
in calibration and glitch model between publications; they must not be
interchanged merely because all are called cleaned data.

There is an additional frequency-band caveat. GWOSC's O2 technical notes
recommend the $16384\,\mathrm{Hz}$ products for studies around
$1700\,\mathrm{Hz}$ and above because the $4096\,\mathrm{Hz}$
anti-alias response corrupts the Nyquist region. The $1875\,\mathrm{Hz}$
diagnostic cannot be promoted to physical high-frequency calibration
without explicit release and filter checks. No frequency cut, data
replacement, gating, PSD regeneration or new inference job was applied.
The frozen calibration bundle remains a legacy mixed-release fixed-PSD
comparison; its historical unresolved-glitch caveat is now qualified by
this independent sample audit, not rewritten in its frozen configuration.

Earlier tables and score snapshots are preserved. The historical
``at least $88\%$'' localization claim above was not reproduced by a
later direct net-score integration, which gave approximately $67\%$
for the $1871$--$1879\,\mathrm{Hz}$ share of the net V1 shift; these
definitions or snapshots must not be combined. Neither those score tests
nor the present release-identification checks validate a detector-consistent
PSD likelihood or a completed GW170817 recovery. A fair estimator
comparison requires PSD estimation and on-source strain from a matched
release with explicit conditioning and a defensible analysis band.

The independent original-input audit saves official API snapshots,
file and sample hashes, aligned spectra, ratios and raw/bandpassed glitch
plots in a fresh timestamped directory. All $16$ hashed source/config/input
files in the frozen calibration bundle were checked without modification.
Sources are the GWOSC GWTC-1 documentation, the GW170817 event-v2/v3
release records, the dedicated GW170817 strain documentation and the O2
technical details; URLs and downloaded-file hashes are recorded in the
audit report. No legacy inference results were replaced.

\subsection{Top-Candidate Cross-Checks}

The same detector- and frequency-resolved normalized-score diagnostic was
applied to the two leading subthreshold candidates, GRB190728271 and
GRB240715239. Unlike GW170817, their discrepancies do not support a common
V1 feature near $1875\,\mathrm{Hz}$. Table~\ref{tab:top-candidate-psd-scores}
uses the sign convention $\ell^{\rm BL}-\ell^{\rm Welch}$; a positive value
therefore indicates a higher normalized fixed-PSD score for BayesLine.

\begin{table}[ht]
\centering
\small
\caption{Detector-resolved normalized fixed-PSD score differences for the
top candidates. The CBC entries use the conditioned maximum-log-likelihood
posterior waveform with the common coalescence phase profiled separately for
each PSD. These are predictive-score diagnostics, not Bayes factors.}
\label{tab:top-candidate-psd-scores}
\begin{tabular}{llrr}
\toprule
\textbf{Event} & \textbf{Detector} & \textbf{Noise-only}
& \textbf{CBC residual}\\
& & \textbf{$\ell^{\rm BL}-\ell^{\rm Welch}$} &
\textbf{$\ell^{\rm BL}-\ell^{\rm Welch}$}\\
\midrule
GRB190728271 & H1 & $+3031.38$ & $+3032.01$\\
GRB190728271 & L1 & $+4583.71$ & $+4581.30$\\
GRB190728271 & Network & $+7615.09$ & $+7613.31$\\
\addlinespace
GRB240715239 & L1 & $+4379.88$ & $+4377.77$\\
GRB240715239 & V1 & $+7616.29$ & $+7622.69$\\
GRB240715239 & Network & $+11996.16$ & $+12000.46$\\
\bottomrule
\end{tabular}
\end{table}

GRB190728271 was observed only in H1 and L1, so a V1 masking test is not
defined for that event. Its score difference is distributed across both
detectors. The largest $1\,\mathrm{Hz}$ contributions occur around
$1461$--$1465\,\mathrm{Hz}$ in H1 and $1022$--$1027\,\mathrm{Hz}$ in L1,
showing that its PSD sensitivity cannot be attributed to the GW170817 V1
feature.

For GRB240715239, dropping V1 leaves the L1-only normalized difference of
$+4379.88$ nats, rather than reversing the sign as it did for GW170817. V1 is
important, contributing $+7616.29$ nats, but L1 also contributes materially.
The targeted V1 interval $1871$--$1879\,\mathrm{Hz}$ contributes only
$+232.12$ nats, or $3.05\%$ of the V1 noise-only difference. Its largest V1
score differences instead lie predominantly between $31$ and
$64\,\mathrm{Hz}$. Thus, although both candidates are sensitive to the fixed
PSD choice, neither has evidence for the specific GW170817-like
$1875\,\mathrm{Hz}$ V1 discrepancy.

As above, detector removal or frequency masking changes the likelihood.
These results localize the fixed-PSD sensitivity, but they do not yield
masked CBC evidences or conditioned-versus-unconditioned Bayes factors.
Those quantities require detector- and mask-consistent likelihood evaluations
and nested-sampling reruns.

The conclusion is therefore methodological. Fixed-PSD evidence differences
should be treated as PSD-sensitivity diagnostics. They motivate inspections
of the affected detector and frequency range, repeat runs using posterior PSD
draws, off-source control studies, and software injections. A definitive
comparison would propagate PSD uncertainty and model non-stationary lines,
window-induced frequency correlations, and transient glitches jointly with the
CBC signal. Finally, removing V1 after the fact permits only an H1/L1
noise-score diagnostic; it does not produce an H1/L1 CBC Bayes factor or
evidence without rerunning the likelihood and nested sampler for that detector
network.

\subsection{Detector- and Frequency-Resolved Likelihood Digest}

For a detector $I$ and frequency bin $k$, define the normalized fixed-PSD
Whittle contribution for a waveform $h$ as
\[
\ell_{I,k}(h;S_n) =
\log\!\left(\frac{2\Delta f}{\pi S_{n,I}(f_k)}\right)
- \frac{2\Delta f\lvert\tilde d_{I,k}-\tilde h_{I,k}\rvert^2}
{S_{n,I}(f_k)}.
\]
The detector score and network score are respectively
\[
\ell_I(h;S_n) = \sum_{k\in[23,2048]~\mathrm{Hz}}\ell_{I,k}(h;S_n),
\qquad
\ell_{\rm net}(h;S_n) = \sum_I\ell_I(h;S_n).
\]
For the noise-only diagnostic, set $h=0$. For the CBC diagnostic, $h$ is a
specified posterior waveform, with a common coalescence phase profiled
separately for each fixed PSD. These are predictive scores for a selected
waveform; they are not nested-sampling evidences and should not replace the
evidence table above.

For direct comparison of the two PSD curves, the detector- and bin-resolved
differences are
\[
\delta\ell_I(h) = \ell_I(h;S_n^{\rm BL})-
\ell_I(h;S_n^{\rm Welch}),
\qquad
\delta\ell_{I,k}(h) =
\ell_{I,k}(h;S_n^{\rm BL})-
\ell_{I,k}(h;S_n^{\rm Welch}),
\]
with $\delta\ell_I=\sum_k\delta\ell_{I,k}$. A positive value means that
BayesLine assigns a higher normalized score; a negative value means that
Welch assigns a higher normalized score. This convention makes the V1
discrepancy transparent: the narrow bin near $1875\,\mathrm{Hz}$ gives a
large negative $\delta\ell_{\mathrm{V1},k}$ when the BayesLine PSD is low
relative to the observed on-source line power.

\begin{table}[ht]
\centering
\small
\caption{Detector-resolved format for the normalized Whittle audit. Each row
is obtained by summing the bin contributions in the stated detector over the
analysis band. The network row is the sum of detector rows. Report the values
from the audit output for the selected event and posterior waveform.}
\begin{tabular}{llrrrr}
\toprule
\textbf{Detector} & \textbf{Waveform} & \textbf{$\ell_I^{\rm Welch}$}
& \textbf{$\ell_I^{\rm BL}$} & \textbf{$\delta\ell_I$}
& \textbf{Interpretation}\\
\midrule
H1 & $h=0$ & \textit{audit value} & \textit{audit value}
& \textit{BL $-$ Welch} & noise-only score\\
L1 & $h=0$ & \textit{audit value} & \textit{audit value}
& \textit{BL $-$ Welch} & noise-only score\\
V1 & $h=0$ & \textit{audit value} & \textit{audit value}
& \textit{BL $-$ Welch} & inspect narrow lines\\
Network & $h=0$ & \textit{sum} & \textit{sum}
& \textit{sum} & detector sum\\
\addlinespace
H1 & $h=h_{\rm prof}$ & \textit{audit value} & \textit{audit value}
& \textit{BL $-$ Welch} & CBC residual score\\
L1 & $h=h_{\rm prof}$ & \textit{audit value} & \textit{audit value}
& \textit{BL $-$ Welch} & CBC residual score\\
V1 & $h=h_{\rm prof}$ & \textit{audit value} & \textit{audit value}
& \textit{BL $-$ Welch} & CBC residual score\\
Network & $h=h_{\rm prof}$ & \textit{sum} & \textit{sum}
& \textit{sum} & detector sum\\
\bottomrule
\end{tabular}
\end{table}

Rather than tabulating every $1/128\,\mathrm{Hz}$ Fourier bin, which would
obscure the result, the per-bin audit lists the bins or narrow bands with the
largest absolute score differences. The confirmed V1 feature can be reported
compactly as follows:
\begin{table}[ht]
\centering
\small
\caption{Largest frequency-resolved contributions to the fixed-PSD score
difference. Values are computed from $\delta\ell_{I,k}$; neighbouring bins
may be grouped only when their summed contribution is reported.}
\begin{tabular}{llrrrl}
\toprule
\textbf{Event} & \textbf{Detector} & \textbf{Frequency/band [Hz]}
& \textbf{$\delta\ell_{I,k}(0)$} & \textbf{$\delta\ell_{I,k}(h_{\rm prof})$}
& \textbf{Feature}\\
\midrule
GRB170817529 & V1 & $1871$--$1879$ &
$\sum_{k\in\mathrm{band}}\delta\ell_{\rm V1,k}(0)
\leq-0.88\lvert\delta\ell_{\rm V1}(0)\rvert$
& \textit{audit value} & narrow feature and leakage skirts\\
\textit{event} & \textit{IFO} & \textit{largest remaining bin/band}
& \textit{audit value} & \textit{audit value} & \textit{diagnostic label}\\
\textit{event} & \textit{IFO} & \textit{largest remaining bin/band}
& \textit{audit value} & \textit{audit value} & \textit{diagnostic label}\\
\bottomrule
\end{tabular}
\end{table}

The two tables should always be read together. A large detector total can be
broad-band, or it can be dominated by one narrow instrumental feature. The
per-bin table distinguishes these cases and prevents a narrow V1 line from
being misdescribed as broad astrophysical residual power. after the Diagnostics and Interpretation Strategy section.
\section{Fixed-PSD Likelihood Audit and Interpretation}

\paragraph{Scope.}
The tables retain historical fixed-PSD audit results. They are not a current
released-map convergence certificate or evidence that BayesLine outperforms
Welch; later input-release qualifications below remain essential.

An additional audit was performed for GRB170817529/GW170817 and
GRB240715239 to make clear which likelihood quantities can be compared across
PSD choices. The sampling likelihood originally used in the pipeline is a
CBC-versus-noise likelihood ratio. For a fixed data segment and a fixed PSD,
this is the appropriate quantity for comparing the conditioned and
unconditioned sky priors: the omitted noise-only term is common to both
models and therefore cancels in $\Delta\log Z$. It is not, however, an
absolute likelihood that can be compared directly between Welch and BayesLine
PSDs.

For the audit, the fixed-PSD one-sided Whittle noise term was evaluated
post-hoc using the same $128\,\mathrm{s}$ Tukey-windowed data and frequency
grid as the sampler:
\[
\log p(d\mid N,S_n) = \sum_k \left[
\log\!\left(\frac{2\Delta f}{\pi S_n(f_k)}\right)
- \frac{2\Delta f\lvert\tilde d_k\rvert^2}{S_n(f_k)}
\right].
\]
This gives the bookkeeping identity
\[
\log Z_{\rm CBC} = \log p(d\mid N,S_n)
+ \log\mathcal{B}_{\rm CBC/N}.
\]
The first term depends on both the PSD determinant and the residual weighting,
so it is not a universal additive normalization. Changing the PSD changes the
relative weights of frequency bins and can consequently change the likelihood
geometry, posterior, and Bayes factor.

\begin{table}[ht]
\centering
\caption{Post-hoc fixed-PSD CBC-versus-noise results. Values are natural
logarithms; uncertainties are nested-sampling bootstrap estimates. The Welch
rows use the established $32\,\mathrm{s}$, $50\%$-overlap median estimator.}
\begin{tabular}{llrrr}
\toprule
\textbf{Event} & \textbf{PSD} & \textbf{$\log\mathcal{B}_{\rm CBC/N}^{\rm cond}$}
& \textbf{$\log\mathcal{B}_{\rm CBC/N}^{\rm uncond}$}
& \textbf{$\Delta\log Z$}\\
\midrule
GRB170817529 & BayesLine & $498.67 \pm 0.31$ & $491.80 \pm 0.34$ & $6.88 \pm 0.46$\\
GRB170817529 & Welch & $522.70 \pm 0.31$ & $517.48 \pm 0.33$ & $5.22 \pm 0.45$\\
GRB240715239 & BayesLine & $8.91 \pm 0.23$ & $5.74 \pm 0.23$ & $3.17 \pm 0.33$\\
GRB240715239 & Welch & $15.46 \pm 0.23$ & $13.54 \pm 0.27$ & $1.92 \pm 0.35$\\
\bottomrule
\end{tabular}
\end{table}

\begin{table}[ht]
\centering
\small
\caption{Post-hoc normalized fixed-PSD likelihood statistics. The normalized
CBC evidence is $\log Z_{\rm CBC}=\log p(d\mid N,S_n)+
\log\mathcal{B}_{\rm CBC/N}$, using the fixed-PSD Whittle convention. Values
are natural logarithms; uncertainties are nested-sampling bootstrap estimates.}
\begin{tabular}{llrrrr}
\toprule
\textbf{Event} & \textbf{PSD/prior} & \textbf{$\log p(d\mid N,S_n)$}
& \textbf{$\log Z_{\rm CBC}$} & \textbf{Mean $\log L$}
& \textbf{$D_{\rm KL}$ (nats)}\\
\midrule
GRB170817529 & BayesLine, conditioned & $76513731.438$ & $76514230.110$
& $76514266.810 \pm 0.086$ & $36.699 \pm 0.294$\\
GRB170817529 & BayesLine, unconditioned & $76513731.438$ & $76514223.233$
& $76514266.781 \pm 0.074$ & $43.547 \pm 0.341$\\
GRB170817529 & Welch, conditioned & $76534785.939$ & $76535308.640$
& $76535345.285 \pm 0.082$ & $36.646 \pm 0.294$\\
GRB170817529 & Welch, unconditioned & $76534785.939$ & $76535303.423$
& $76535345.359 \pm 0.079$ & $41.936 \pm 0.329$\\
\addlinespace
GRB240715239 & BayesLine, conditioned & $51677741.063$ & $51677749.975$
& $51677771.095 \pm 0.260$ & $21.120 \pm 0.268$\\
GRB240715239 & BayesLine, unconditioned & $51677741.063$ & $51677746.807$
& $51677766.637 \pm 0.136$ & $19.830 \pm 0.189$\\
GRB240715239 & Welch, conditioned & $51658062.380$ & $51658077.842$
& $51658098.695 \pm 0.112$ & $20.853 \pm 0.216$\\
GRB240715239 & Welch, unconditioned & $51658062.380$ & $51658075.921$
& $51658103.335 \pm 0.172$ & $27.413 \pm 0.250$\\
\bottomrule
\end{tabular}
\end{table}

These results show genuine sensitivity to the fixed PSD, rather than evidence
for one missing scalar normalization. In particular, GRB240715239 retains a
positive conditioned-versus-unconditioned comparison under both PSD choices,
but its magnitude changes. GW170817 remains decisively CBC-supported under
both choices and remains a useful control. The very large absolute values of
$\log p(d\mid N,S_n)$ and $\log Z_{\rm CBC}$ should not be interpreted as
an astrophysical ranking of PSD methods: they are dominated by the many-bin
noise likelihood and depend on the chosen fixed-PSD convention.

\subsection{Why the Welch--BayesLine Residual Difference Is Not ``False Physics''}

\textbf{Historical interpretation, qualified by the original-input audit
of 1 November 2026 below.} The PSD curves were subsequently shown to use
different published strain releases. The discussion of estimator and window
effects here must therefore not be treated as an identical-data comparison.

The large residual seen in the Welch-minus-BayesLine score is not generated by
Welch as an astrophysical feature. It is a consequence of the two fixed PSD
curves assigning very different noise power to a narrow instrumental feature,
most visibly in V1 near $1875\,\mathrm{Hz}$ for GW170817. Inspection of the
Welch PSD shows the broad, approximately symmetric skirts expected from
spectral leakage around a narrow feature. Across those skirts, the fixed Welch
and BayesLine PSDs differ by orders of magnitude in some bins. Bins with low
observed Fourier power are consequently assigned very different Whittle
weights, $1/S_n(f_k)$, and can make a large cumulative contribution to the
residual score. Expressed as a Welch-minus-BayesLine score, this produces a
large positive V1 contribution.

This does \emph{not} establish that Welch is physically correct or that
BayesLine is physically false. It identifies a calibration and modelling
issue: a posterior-median, off-source BayesLine PSD can under-represent the
power of a narrow line in the particular on-source realization, while a
median Welch estimate can absorb that line power more conservatively. A
held-out calibration test supports this interpretation. BayesLine is close to
the exponential-noise reference in the broad H1/L1/V1 distributions, while
the $32\,\mathrm{s}$ Welch estimate has a heavier upper tail, consistent with
localized excess power relative to its fixed PSD. Neither observation alone
selects a complete physical noise model for the on-source data.

Welch has no explicit transient-glitch model analogous to the wavelet
component available in BayesWave alongside its BayesLine PSD model. Median
averaging can reduce the influence of outlying periodograms, but it does not
identify, reconstruct, or subtract a glitch. This distinction must be kept
separate from spectral leakage caused by the finite-data window.

\subsection{Confirmation and Masking Tests}

The origin of the discrepancy was tested in three complementary ways. First,
the detector-resolved normalized score identified V1 as the only detector
with a large positive Welch-minus-BayesLine contribution. Second, a
frequency-resolved sum showed that at least $88\%$ of that V1 discrepancy is
concentrated in the narrow interval $1871$--$1879\,\mathrm{Hz}$, rather than
being distributed across the astrophysically relevant inspiral band. Third,
BayesWave/BayesLine was rerun with its glitch-wavelet component disabled. The no-glitch PSD did
not show a material drift from the standard glitch-modelled BayesLine PSD in
this region. The recorded comparison therefore showed no indication that
allocating transient power to glitch wavelets was responsible for the
discrepancy with Welch. It supports a narrow-feature/PSD-estimation or
window-response explanation rather than a difference in glitch treatment;
it does not, by itself, prove that no transient was present. This V1
high-frequency control is distinct from the known on-source L1 glitch in
GW170817, which an off-source PSD analysis does not remove.

\begin{table}[ht]
\centering
\small
\caption{GW170817 detector-resolved normalized fixed-PSD score difference,
$\ell^{\rm Welch}-\ell^{\rm BL}$, before any masking. These are absolute
fixed-PSD likelihood scores, not CBC Bayes factors.}
\begin{tabular}{lr}
\toprule
\textbf{Detector} & \textbf{$\ell^{\rm Welch}-\ell^{\rm BL}$ (nats)}\\
\midrule
H1 & $-1941.71$\\
L1 & $-4770.43$\\
V1 & $+27766.64$\\
\midrule
Network & $+21054.50$\\
\bottomrule
\end{tabular}
\end{table}

The masking tests are decisive as a diagnostic of the source of the absolute
score difference:
\begin{itemize}
\item Masking only $1871$--$1879\,\mathrm{Hz}$ removes at least $88\%$ of the
V1 positive contribution. This localizes the effect to the narrow spectral
feature and its leakage structure.
\item Excluding V1 entirely changes the network Welch-minus-BayesLine
normalized log likelihood from $+21054.50$ to $-6712.14$ nats. The sign flip
shows that V1, rather than H1 or L1, controls the original positive network
difference.
\end{itemize}

Neither operation is a replacement for a detector-consistent Bayesian
analysis. Frequency masking changes the likelihood and must be applied to
both the CBC and noise hypotheses before recomputing an evidence. Similarly,
the V1-excluded number is an H1/L1 noise-score diagnostic; an H1/L1 CBC Bayes
factor or evidence requires rerunning the likelihood and nested sampler with
V1 absent.

\subsection{Window-Response and Line-Evolution Test}

A follow-up on 30 October 2026 tested the finite-window response separately
from evolution of the dominant V1 line. The $1024\,\mathrm{s}$ off-source
context was divided into eight non-overlapping $128\,\mathrm{s}$ blocks.
Only the first four blocks were used to train the predictions; the remaining
four blocks and the on-source interval were withheld. An independent fit
in each block was retained as a diagnostic control, not as a held-out
prediction. These are narrow-band tests, not new CBC evidence calculations.

The dominant line was fitted near $1875.39202\,\mathrm{Hz}$, with an
all-block fitted-frequency range of only $3.8\times10^{-5}\,\mathrm{Hz}$.
Its peak strain amplitude nevertheless decreased from approximately
$9.46\times10^{-21}$ in the earliest block to $4.28\times10^{-21}$
on-source, corresponding to about an $80\%$ decline in line power.
A log-linear amplitude trend fitted exclusively to the four training blocks
had an amplitude e-folding time of approximately $1389\,\mathrm{s}$.
It predicted the later fitted line powers within about $6\%$, and the
on-source fitted line power within about $9\%$. This demonstrates amplitude
evolution in the cached strain; it does not identify the physical cause or
establish a transient-glitch classification.

The forward prediction used a sinusoid with that training-only frequency
and amplitude trend, including its within-block envelope, plus an estimated
local continuum. Its exact phase-averaged response was computed through
the symmetric Tukey window used by the on-source FFT. This did not convolve
an already-windowed PSD a second time. Synthetic-tone recovery,
window-weighted power conservation, and agreement with the pipeline's
on-source FFT periodogram were checked.

For the actual on-source Tukey window, the observed-to-predicted integrated
power ratios were $0.913$ in the core and $1.044$ in the skirts. The core
was defined within $0.25\,\mathrm{Hz}$ of the training frequency; the skirts
were $0.25$--$4\,\mathrm{Hz}$ away. Across the four held-out off-source
blocks these ratios ranged from $0.938$ to $1.004$ in the core and from
$1.022$ to $1.142$ in the skirts. In contrast, the stationary-amplitude
prediction gave an on-source core ratio of only $0.269$: matching the window
alone did not account for the change in line strength.

Within the local cache, this establishes a finite-window response and
nonstationary line amplitude. It does not by itself establish their share
of the cross-release Welch--BayesLine discrepancy. Agreement in
integrated power is not a complete calibration: the on-source mean
per-bin skirt-power ratio was still about $1.69$, versus a median of
$0.95$, indicating remaining localized excess. A Hann-window diagnostic
exposed additional narrow structures, leaving skirt power about $14$ times
the single-line-plus-continuum prediction. Parameter uncertainty,
additional lines, and frequency-bin covariance remain unmodelled. Local
core-conditioned fits are explanatory controls, and the reported diagonal
noise-score shifts are not model-selection Bayes factors. The BayesLine PSD
was displayed only as an on-source reference, not as an independently
trained held-out model. Production windows, PSD settings, and queued jobs
were unchanged.

\subsection{Matched No-Glitch BayesLine Control}

A diagnostic on 30 October 2026 matched the nominal Welch settings to the
archived \texttt{noiseOnly}, \texttt{noClean}, \texttt{noWaveletPrior}
BayesLine run. Its analysis interval was GPS
$[1187008610.4,1187008738.4]$, with $4096\,\mathrm{Hz}$ sampling,
a $128\,\mathrm{s}$ FFT, $\Delta f=0.0078125\,\mathrm{Hz}$, and the
$23$--$2048\,\mathrm{Hz}$ band. A $0.4\,\mathrm{s}$ edge taper,
corresponding to Tukey $\alpha=0.00625$, reproduces the logged window
correction to its printed precision. The separate configuration option
\texttt{window=1.0} was not interpreted as a Tukey alpha.

One complete segment produces a single periodogram, not a variance-reduced
Welch average. Its normalization was checked against a direct FFT and
window-weighted Parseval identity. Raw posterior PSDs were converted by
$2/128$ and checked against the archived physical-units export; the same
conversion was applied to the archived starting PSD. BayesWavePost omits
Nyquist and rounds frequency labels, so ordered rows were checked against
the exact starting-PSD grid instead of interpolating the rounded labels.

The nominally matched local strain still disagreed strongly with the
BayesLine PSD. However, the archived no-glitch Fourier residual export
revealed a more important distinction: it was not numerically equivalent
to the local strain. Integrated data-to-BayesLine power ratios were
\[
\begin{array}{lrr}
 & \text{Local cache} & \text{Archived Fourier export}\\
\text{Core: within }0.25\,\mathrm{Hz} & 1312.61 & 0.693\\
\text{Skirts: }0.25\text{--}4\,\mathrm{Hz} & 10.993 & 0.900\\
\text{Background: }5\text{--}8\,\mathrm{Hz} & 1.036 & 0.971
\end{array}
\]
around the fitted $1875.3920\,\mathrm{Hz}$ line. The archived export had
about $1900$ times less integrated core power than the local cache.
Symmetric versus periodic Tukey windows changed local core power by only
about $3.1\times10^{-7}$ fractionally. The nearest local sample starts
$0.4004$ samples earlier than the requested GPS time; adjacent-sample
sensitivity controls were also recorded.

The executed run flags disable glitch, signal, and CBC components, making
its Fourier residual a useful processed-data control. Nevertheless, it is
an exported residual product, not an independently verified copy of the
unconditioned frame strain. These results therefore locate an unresolved
data-provenance, conditioning, or export-consistency issue; they do not
identify its cause, demonstrate deliberate line removal, or establish that
BayesLine intrinsically misses the line. The archive-only core and skirt
ratios are of order one, but this is an in-sample diagnostic, not noise-model
calibration. The earlier conclusion that matching nominal settings alone
tests estimator differences must be qualified by this processed-data check.

The full $1024\,\mathrm{s}$ initialization context was not reproduced:
its first $2.075\,\mathrm{s}$ are absent from the local cache. No new
BayesWave or CBC inference jobs were launched, and production settings were
unchanged. Numerical checks, source hashes, arrays, and plots are saved by
the \texttt{--compare-no-glitch} diagnostic in its separate
\texttt{matched\_no\_glitch} output directory.

\subsection{Original-Input Release Identification and L1 Glitch Status}

On 1 November 2026 the actual frame input was checked rather than relying
on the post-processing export or nominal estimator settings. The archived
H1/L1/V1 frame-cache records name the $4096\,\mathrm{Hz}$ GWTC-1
GW170817 event-v3 files beginning at GPS $1187006835$, with duration
$4096\,\mathrm{s}$. SHA256 hashes of all three local frames matched their
current official GWOSC download bytes exactly. Duplicate local frame
copies also matched. This authenticates the available original input;
it is not a recovery of the historical download log.

The local sampler cache was separately compared with the official
O1\_O2-Preliminary GW170817 event-v2 cleaned HDF5 products beginning at
GPS $1187007040$. \emph{Every one of the $4{,}784{,}128$ aligned samples
per detector matched exactly}, over the complete cached
$1168\,\mathrm{s}$ interval
$[1187007716.4746094,1187008884.4746094)$. All H1/L1/V1 cache copies
checked against the local historical archive also had identical file
hashes. Thus the cached strain is identified as the published v2 release,
not as an unspecified O2 bulk product. The collection code's
\texttt{fetch\_open\_data(version=2)} is consistent with this finding,
but the identification rests on sample equality, not the version argument.

For the nominal $128\,\mathrm{s}$ no-glitch control, the first frame
sample at or after the requested start is GPS $1187008610.4001465$.
Using the same aligned local samples, symmetric Tukey taper and
periodogram normalization gives original-frame/local-cache integrated
V1 power ratios
\[
\begin{array}{lr}
\text{Core: within }0.25\,\mathrm{Hz}\text{ of }1875.3920\,\mathrm{Hz}
 & 0.00053034\\
\text{Skirts: }0.25\text{--}4\,\mathrm{Hz} & 0.0822164\\
\text{Broadband: }100\text{--}1500\,\mathrm{Hz} & 0.936858.
\end{array}
\]
The original frame already has approximately $1886$ times less core power.
The archived V1 Fourier export closely reproduces that frame: the relative
complex residual is approximately $0.00119$, with fitted complex scale
$0.99999908+0.00000081i$. Its integrated power divided by the physical
frame periodogram is approximately $0.996093$, consistent with the
window-energy factor $\langle w^2\rangle$ omitted from the export's
convention. A one-sample difference between nearest-sample and
first-at-or-after-start extraction materially affects the complex phases,
but not the orders-of-magnitude power distinction.

\textbf{The large V1 discrepancy is therefore present in the input strain
before BayesLine fitting.} It is not evidence that BayesLine intrinsically
misses the strong line seen in a different release. The successful complex
match also rules out a wholly independent simulated-noise realization as
an adequate explanation for this export, despite the simulated-channel
labels in the BayesWavePost log. The exact remaining conditioning,
calibration and downsampling differences have not been isolated, and the
precise historical post-processing overwrite sequence is not inferred.

The known strong GW170817 glitch is a separate L1 issue. GWOSC documents
its time as GPS $1187008881.389$, about $1.1\,\mathrm{s}$ before merger,
and a brief digital-to-analog converter saturation of less than
$5\,\mathrm{ms}$. GWTC-1 event-v3 strain retains the glitch; the earlier
v2 release is explicitly documented as glitch-subtracted. Since the full
cached L1 samples exactly match v2, \emph{the local sampler cache already
uses the published glitch-subtracted L1 data}. A common diagnostic
$30$--$1000\,\mathrm{Hz}$ zero-phase bandpass shows the transient in the
original v3 frame and its suppression in the cache. This plot is a display
control, not a reconstruction-validation test or a new likelihood.
Unfiltered raw maxima alone are misleading because the v2 H1/L1 products
have documented excess noise below $20\,\mathrm{Hz}$. The v2 H1
last-$150\,\mathrm{s}$ cleaning-window warning concerns approximately
GPS $1187008938$ onward, outside the cached interval checked here.

The off-source BayesWave control ends at GPS $1187008738.4$, before the
L1 glitch. Off-source PSD cleaning neither explains the cross-detector
V1 difference nor establishes on-source glitch subtraction. In this case,
the cache's subtraction status is established independently by its release
content. Special glitch-subtracted products described by GWOSC also differ
in calibration and glitch model between publications; they must not be
interchanged merely because all are called cleaned data.

There is an additional frequency-band caveat. GWOSC's O2 technical notes
recommend the $16384\,\mathrm{Hz}$ products for studies around
$1700\,\mathrm{Hz}$ and above because the $4096\,\mathrm{Hz}$
anti-alias response corrupts the Nyquist region. The $1875\,\mathrm{Hz}$
diagnostic cannot be promoted to physical high-frequency calibration
without explicit release and filter checks. No frequency cut, data
replacement, gating, PSD regeneration or new inference job was applied.
The frozen calibration bundle remains a legacy mixed-release fixed-PSD
comparison; its historical unresolved-glitch caveat is now qualified by
this independent sample audit, not rewritten in its frozen configuration.

Earlier tables and score snapshots are preserved. The historical
``at least $88\%$'' localization claim above was not reproduced by a
later direct net-score integration, which gave approximately $67\%$
for the $1871$--$1879\,\mathrm{Hz}$ share of the net V1 shift; these
definitions or snapshots must not be combined. Neither those score tests
nor the present release-identification checks validate a detector-consistent
PSD likelihood or a completed GW170817 recovery. A fair estimator
comparison requires PSD estimation and on-source strain from a matched
release with explicit conditioning and a defensible analysis band.

The independent original-input audit saves official API snapshots,
file and sample hashes, aligned spectra, ratios and raw/bandpassed glitch
plots in a fresh timestamped directory. All $16$ hashed source/config/input
files in the frozen calibration bundle were checked without modification.
Sources are the GWOSC GWTC-1 documentation, the GW170817 event-v2/v3
release records, the dedicated GW170817 strain documentation and the O2
technical details; URLs and downloaded-file hashes are recorded in the
audit report. No legacy inference results were replaced.

\subsection{Top-Candidate Cross-Checks}

The same detector- and frequency-resolved normalized-score diagnostic was
applied to the two leading subthreshold candidates, GRB190728271 and
GRB240715239. Unlike GW170817, their discrepancies do not support a common
V1 feature near $1875\,\mathrm{Hz}$. Table~\ref{tab:top-candidate-psd-scores}
uses the sign convention $\ell^{\rm BL}-\ell^{\rm Welch}$; a positive value
therefore indicates a higher normalized fixed-PSD score for BayesLine.

\begin{table}[ht]
\centering
\small
\caption{Detector-resolved normalized fixed-PSD score differences for the
top candidates. The CBC entries use the conditioned maximum-log-likelihood
posterior waveform with the common coalescence phase profiled separately for
each PSD. These are predictive-score diagnostics, not Bayes factors.}
\label{tab:top-candidate-psd-scores}
\begin{tabular}{llrr}
\toprule
\textbf{Event} & \textbf{Detector} & \textbf{Noise-only}
& \textbf{CBC residual}\\
& & \textbf{$\ell^{\rm BL}-\ell^{\rm Welch}$} &
\textbf{$\ell^{\rm BL}-\ell^{\rm Welch}$}\\
\midrule
GRB190728271 & H1 & $+3031.38$ & $+3032.01$\\
GRB190728271 & L1 & $+4583.71$ & $+4581.30$\\
GRB190728271 & Network & $+7615.09$ & $+7613.31$\\
\addlinespace
GRB240715239 & L1 & $+4379.88$ & $+4377.77$\\
GRB240715239 & V1 & $+7616.29$ & $+7622.69$\\
GRB240715239 & Network & $+11996.16$ & $+12000.46$\\
\bottomrule
\end{tabular}
\end{table}

GRB190728271 was observed only in H1 and L1, so a V1 masking test is not
defined for that event. Its score difference is distributed across both
detectors. The largest $1\,\mathrm{Hz}$ contributions occur around
$1461$--$1465\,\mathrm{Hz}$ in H1 and $1022$--$1027\,\mathrm{Hz}$ in L1,
showing that its PSD sensitivity cannot be attributed to the GW170817 V1
feature.

For GRB240715239, dropping V1 leaves the L1-only normalized difference of
$+4379.88$ nats, rather than reversing the sign as it did for GW170817. V1 is
important, contributing $+7616.29$ nats, but L1 also contributes materially.
The targeted V1 interval $1871$--$1879\,\mathrm{Hz}$ contributes only
$+232.12$ nats, or $3.05\%$ of the V1 noise-only difference. Its largest V1
score differences instead lie predominantly between $31$ and
$64\,\mathrm{Hz}$. Thus, although both candidates are sensitive to the fixed
PSD choice, neither has evidence for the specific GW170817-like
$1875\,\mathrm{Hz}$ V1 discrepancy.

As above, detector removal or frequency masking changes the likelihood.
These results localize the fixed-PSD sensitivity, but they do not yield
masked CBC evidences or conditioned-versus-unconditioned Bayes factors.
Those quantities require detector- and mask-consistent likelihood evaluations
and nested-sampling reruns.

The conclusion is therefore methodological. Fixed-PSD evidence differences
should be treated as PSD-sensitivity diagnostics. They motivate inspections
of the affected detector and frequency range, repeat runs using posterior PSD
draws, off-source control studies, and software injections. A definitive
comparison would propagate PSD uncertainty and model non-stationary lines,
window-induced frequency correlations, and transient glitches jointly with the
CBC signal. Finally, removing V1 after the fact permits only an H1/L1
noise-score diagnostic; it does not produce an H1/L1 CBC Bayes factor or
evidence without rerunning the likelihood and nested sampler for that detector
network.

\subsection{Detector- and Frequency-Resolved Likelihood Digest}

For a detector $I$ and frequency bin $k$, define the normalized fixed-PSD
Whittle contribution for a waveform $h$ as
\[
\ell_{I,k}(h;S_n) =
\log\!\left(\frac{2\Delta f}{\pi S_{n,I}(f_k)}\right)
- \frac{2\Delta f\lvert\tilde d_{I,k}-\tilde h_{I,k}\rvert^2}
{S_{n,I}(f_k)}.
\]
The detector score and network score are respectively
\[
\ell_I(h;S_n) = \sum_{k\in[23,2048]~\mathrm{Hz}}\ell_{I,k}(h;S_n),
\qquad
\ell_{\rm net}(h;S_n) = \sum_I\ell_I(h;S_n).
\]
For the noise-only diagnostic, set $h=0$. For the CBC diagnostic, $h$ is a
specified posterior waveform, with a common coalescence phase profiled
separately for each fixed PSD. These are predictive scores for a selected
waveform; they are not nested-sampling evidences and should not replace the
evidence table above.

For direct comparison of the two PSD curves, the detector- and bin-resolved
differences are
\[
\delta\ell_I(h) = \ell_I(h;S_n^{\rm BL})-
\ell_I(h;S_n^{\rm Welch}),
\qquad
\delta\ell_{I,k}(h) =
\ell_{I,k}(h;S_n^{\rm BL})-
\ell_{I,k}(h;S_n^{\rm Welch}),
\]
with $\delta\ell_I=\sum_k\delta\ell_{I,k}$. A positive value means that
BayesLine assigns a higher normalized score; a negative value means that
Welch assigns a higher normalized score. This convention makes the V1
discrepancy transparent: the narrow bin near $1875\,\mathrm{Hz}$ gives a
large negative $\delta\ell_{\mathrm{V1},k}$ when the BayesLine PSD is low
relative to the observed on-source line power.

\begin{table}[ht]
\centering
\small
\caption{Detector-resolved format for the normalized Whittle audit. Each row
is obtained by summing the bin contributions in the stated detector over the
analysis band. The network row is the sum of detector rows. Report the values
from the audit output for the selected event and posterior waveform.}
\begin{tabular}{llrrrr}
\toprule
\textbf{Detector} & \textbf{Waveform} & \textbf{$\ell_I^{\rm Welch}$}
& \textbf{$\ell_I^{\rm BL}$} & \textbf{$\delta\ell_I$}
& \textbf{Interpretation}\\
\midrule
H1 & $h=0$ & \textit{audit value} & \textit{audit value}
& \textit{BL $-$ Welch} & noise-only score\\
L1 & $h=0$ & \textit{audit value} & \textit{audit value}
& \textit{BL $-$ Welch} & noise-only score\\
V1 & $h=0$ & \textit{audit value} & \textit{audit value}
& \textit{BL $-$ Welch} & inspect narrow lines\\
Network & $h=0$ & \textit{sum} & \textit{sum}
& \textit{sum} & detector sum\\
\addlinespace
H1 & $h=h_{\rm prof}$ & \textit{audit value} & \textit{audit value}
& \textit{BL $-$ Welch} & CBC residual score\\
L1 & $h=h_{\rm prof}$ & \textit{audit value} & \textit{audit value}
& \textit{BL $-$ Welch} & CBC residual score\\
V1 & $h=h_{\rm prof}$ & \textit{audit value} & \textit{audit value}
& \textit{BL $-$ Welch} & CBC residual score\\
Network & $h=h_{\rm prof}$ & \textit{sum} & \textit{sum}
& \textit{sum} & detector sum\\
\bottomrule
\end{tabular}
\end{table}

Rather than tabulating every $1/128\,\mathrm{Hz}$ Fourier bin, which would
obscure the result, the per-bin audit lists the bins or narrow bands with the
largest absolute score differences. The confirmed V1 feature can be reported
compactly as follows:
\begin{table}[ht]
\centering
\small
\caption{Largest frequency-resolved contributions to the fixed-PSD score
difference. Values are computed from $\delta\ell_{I,k}$; neighbouring bins
may be grouped only when their summed contribution is reported.}
\begin{tabular}{llrrrl}
\toprule
\textbf{Event} & \textbf{Detector} & \textbf{Frequency/band [Hz]}
& \textbf{$\delta\ell_{I,k}(0)$} & \textbf{$\delta\ell_{I,k}(h_{\rm prof})$}
& \textbf{Feature}\\
\midrule
GRB170817529 & V1 & $1871$--$1879$ &
$\sum_{k\in\mathrm{band}}\delta\ell_{\rm V1,k}(0)
\leq-0.88\lvert\delta\ell_{\rm V1}(0)\rvert$
& \textit{audit value} & narrow feature and leakage skirts\\
\textit{event} & \textit{IFO} & \textit{largest remaining bin/band}
& \textit{audit value} & \textit{audit value} & \textit{diagnostic label}\\
\textit{event} & \textit{IFO} & \textit{largest remaining bin/band}
& \textit{audit value} & \textit{audit value} & \textit{diagnostic label}\\
\bottomrule
\end{tabular}
\end{table}

The two tables should always be read together. A large detector total can be
broad-band, or it can be dominated by one narrow instrumental feature. The
per-bin table distinguishes these cases and prevents a narrow V1 line from
being misdescribed as broad astrophysical residual power. after the Diagnostics and Interpretation Strategy section.
\section{Fixed-PSD Likelihood Audit and Interpretation}

\paragraph{Scope.}
The tables retain historical fixed-PSD audit results. They are not a current
released-map convergence certificate or evidence that BayesLine outperforms
Welch; later input-release qualifications below remain essential.

An additional audit was performed for GRB170817529/GW170817 and
GRB240715239 to make clear which likelihood quantities can be compared across
PSD choices. The sampling likelihood originally used in the pipeline is a
CBC-versus-noise likelihood ratio. For a fixed data segment and a fixed PSD,
this is the appropriate quantity for comparing the conditioned and
unconditioned sky priors: the omitted noise-only term is common to both
models and therefore cancels in $\Delta\log Z$. It is not, however, an
absolute likelihood that can be compared directly between Welch and BayesLine
PSDs.

For the audit, the fixed-PSD one-sided Whittle noise term was evaluated
post-hoc using the same $128\,\mathrm{s}$ Tukey-windowed data and frequency
grid as the sampler:
\[
\log p(d\mid N,S_n) = \sum_k \left[
\log\!\left(\frac{2\Delta f}{\pi S_n(f_k)}\right)
- \frac{2\Delta f\lvert\tilde d_k\rvert^2}{S_n(f_k)}
\right].
\]
This gives the bookkeeping identity
\[
\log Z_{\rm CBC} = \log p(d\mid N,S_n)
+ \log\mathcal{B}_{\rm CBC/N}.
\]
The first term depends on both the PSD determinant and the residual weighting,
so it is not a universal additive normalization. Changing the PSD changes the
relative weights of frequency bins and can consequently change the likelihood
geometry, posterior, and Bayes factor.

\begin{table}[ht]
\centering
\caption{Post-hoc fixed-PSD CBC-versus-noise results. Values are natural
logarithms; uncertainties are nested-sampling bootstrap estimates. The Welch
rows use the established $32\,\mathrm{s}$, $50\%$-overlap median estimator.}
\begin{tabular}{llrrr}
\toprule
\textbf{Event} & \textbf{PSD} & \textbf{$\log\mathcal{B}_{\rm CBC/N}^{\rm cond}$}
& \textbf{$\log\mathcal{B}_{\rm CBC/N}^{\rm uncond}$}
& \textbf{$\Delta\log Z$}\\
\midrule
GRB170817529 & BayesLine & $498.67 \pm 0.31$ & $491.80 \pm 0.34$ & $6.88 \pm 0.46$\\
GRB170817529 & Welch & $522.70 \pm 0.31$ & $517.48 \pm 0.33$ & $5.22 \pm 0.45$\\
GRB240715239 & BayesLine & $8.91 \pm 0.23$ & $5.74 \pm 0.23$ & $3.17 \pm 0.33$\\
GRB240715239 & Welch & $15.46 \pm 0.23$ & $13.54 \pm 0.27$ & $1.92 \pm 0.35$\\
\bottomrule
\end{tabular}
\end{table}

\begin{table}[ht]
\centering
\small
\caption{Post-hoc normalized fixed-PSD likelihood statistics. The normalized
CBC evidence is $\log Z_{\rm CBC}=\log p(d\mid N,S_n)+
\log\mathcal{B}_{\rm CBC/N}$, using the fixed-PSD Whittle convention. Values
are natural logarithms; uncertainties are nested-sampling bootstrap estimates.}
\begin{tabular}{llrrrr}
\toprule
\textbf{Event} & \textbf{PSD/prior} & \textbf{$\log p(d\mid N,S_n)$}
& \textbf{$\log Z_{\rm CBC}$} & \textbf{Mean $\log L$}
& \textbf{$D_{\rm KL}$ (nats)}\\
\midrule
GRB170817529 & BayesLine, conditioned & $76513731.438$ & $76514230.110$
& $76514266.810 \pm 0.086$ & $36.699 \pm 0.294$\\
GRB170817529 & BayesLine, unconditioned & $76513731.438$ & $76514223.233$
& $76514266.781 \pm 0.074$ & $43.547 \pm 0.341$\\
GRB170817529 & Welch, conditioned & $76534785.939$ & $76535308.640$
& $76535345.285 \pm 0.082$ & $36.646 \pm 0.294$\\
GRB170817529 & Welch, unconditioned & $76534785.939$ & $76535303.423$
& $76535345.359 \pm 0.079$ & $41.936 \pm 0.329$\\
\addlinespace
GRB240715239 & BayesLine, conditioned & $51677741.063$ & $51677749.975$
& $51677771.095 \pm 0.260$ & $21.120 \pm 0.268$\\
GRB240715239 & BayesLine, unconditioned & $51677741.063$ & $51677746.807$
& $51677766.637 \pm 0.136$ & $19.830 \pm 0.189$\\
GRB240715239 & Welch, conditioned & $51658062.380$ & $51658077.842$
& $51658098.695 \pm 0.112$ & $20.853 \pm 0.216$\\
GRB240715239 & Welch, unconditioned & $51658062.380$ & $51658075.921$
& $51658103.335 \pm 0.172$ & $27.413 \pm 0.250$\\
\bottomrule
\end{tabular}
\end{table}

These results show genuine sensitivity to the fixed PSD, rather than evidence
for one missing scalar normalization. In particular, GRB240715239 retains a
positive conditioned-versus-unconditioned comparison under both PSD choices,
but its magnitude changes. GW170817 remains decisively CBC-supported under
both choices and remains a useful control. The very large absolute values of
$\log p(d\mid N,S_n)$ and $\log Z_{\rm CBC}$ should not be interpreted as
an astrophysical ranking of PSD methods: they are dominated by the many-bin
noise likelihood and depend on the chosen fixed-PSD convention.

\subsection{Why the Welch--BayesLine Residual Difference Is Not ``False Physics''}

\textbf{Historical interpretation, qualified by the original-input audit
of 1 November 2026 below.} The PSD curves were subsequently shown to use
different published strain releases. The discussion of estimator and window
effects here must therefore not be treated as an identical-data comparison.

The large residual seen in the Welch-minus-BayesLine score is not generated by
Welch as an astrophysical feature. It is a consequence of the two fixed PSD
curves assigning very different noise power to a narrow instrumental feature,
most visibly in V1 near $1875\,\mathrm{Hz}$ for GW170817. Inspection of the
Welch PSD shows the broad, approximately symmetric skirts expected from
spectral leakage around a narrow feature. Across those skirts, the fixed Welch
and BayesLine PSDs differ by orders of magnitude in some bins. Bins with low
observed Fourier power are consequently assigned very different Whittle
weights, $1/S_n(f_k)$, and can make a large cumulative contribution to the
residual score. Expressed as a Welch-minus-BayesLine score, this produces a
large positive V1 contribution.

This does \emph{not} establish that Welch is physically correct or that
BayesLine is physically false. It identifies a calibration and modelling
issue: a posterior-median, off-source BayesLine PSD can under-represent the
power of a narrow line in the particular on-source realization, while a
median Welch estimate can absorb that line power more conservatively. A
held-out calibration test supports this interpretation. BayesLine is close to
the exponential-noise reference in the broad H1/L1/V1 distributions, while
the $32\,\mathrm{s}$ Welch estimate has a heavier upper tail, consistent with
localized excess power relative to its fixed PSD. Neither observation alone
selects a complete physical noise model for the on-source data.

Welch has no explicit transient-glitch model analogous to the wavelet
component available in BayesWave alongside its BayesLine PSD model. Median
averaging can reduce the influence of outlying periodograms, but it does not
identify, reconstruct, or subtract a glitch. This distinction must be kept
separate from spectral leakage caused by the finite-data window.

\subsection{Confirmation and Masking Tests}

The origin of the discrepancy was tested in three complementary ways. First,
the detector-resolved normalized score identified V1 as the only detector
with a large positive Welch-minus-BayesLine contribution. Second, a
frequency-resolved sum showed that at least $88\%$ of that V1 discrepancy is
concentrated in the narrow interval $1871$--$1879\,\mathrm{Hz}$, rather than
being distributed across the astrophysically relevant inspiral band. Third,
BayesWave/BayesLine was rerun with its glitch-wavelet component disabled. The no-glitch PSD did
not show a material drift from the standard glitch-modelled BayesLine PSD in
this region. The recorded comparison therefore showed no indication that
allocating transient power to glitch wavelets was responsible for the
discrepancy with Welch. It supports a narrow-feature/PSD-estimation or
window-response explanation rather than a difference in glitch treatment;
it does not, by itself, prove that no transient was present. This V1
high-frequency control is distinct from the known on-source L1 glitch in
GW170817, which an off-source PSD analysis does not remove.

\begin{table}[ht]
\centering
\small
\caption{GW170817 detector-resolved normalized fixed-PSD score difference,
$\ell^{\rm Welch}-\ell^{\rm BL}$, before any masking. These are absolute
fixed-PSD likelihood scores, not CBC Bayes factors.}
\begin{tabular}{lr}
\toprule
\textbf{Detector} & \textbf{$\ell^{\rm Welch}-\ell^{\rm BL}$ (nats)}\\
\midrule
H1 & $-1941.71$\\
L1 & $-4770.43$\\
V1 & $+27766.64$\\
\midrule
Network & $+21054.50$\\
\bottomrule
\end{tabular}
\end{table}

The masking tests are decisive as a diagnostic of the source of the absolute
score difference:
\begin{itemize}
\item Masking only $1871$--$1879\,\mathrm{Hz}$ removes at least $88\%$ of the
V1 positive contribution. This localizes the effect to the narrow spectral
feature and its leakage structure.
\item Excluding V1 entirely changes the network Welch-minus-BayesLine
normalized log likelihood from $+21054.50$ to $-6712.14$ nats. The sign flip
shows that V1, rather than H1 or L1, controls the original positive network
difference.
\end{itemize}

Neither operation is a replacement for a detector-consistent Bayesian
analysis. Frequency masking changes the likelihood and must be applied to
both the CBC and noise hypotheses before recomputing an evidence. Similarly,
the V1-excluded number is an H1/L1 noise-score diagnostic; an H1/L1 CBC Bayes
factor or evidence requires rerunning the likelihood and nested sampler with
V1 absent.

\subsection{Window-Response and Line-Evolution Test}

A follow-up on 30 October 2026 tested the finite-window response separately
from evolution of the dominant V1 line. The $1024\,\mathrm{s}$ off-source
context was divided into eight non-overlapping $128\,\mathrm{s}$ blocks.
Only the first four blocks were used to train the predictions; the remaining
four blocks and the on-source interval were withheld. An independent fit
in each block was retained as a diagnostic control, not as a held-out
prediction. These are narrow-band tests, not new CBC evidence calculations.

The dominant line was fitted near $1875.39202\,\mathrm{Hz}$, with an
all-block fitted-frequency range of only $3.8\times10^{-5}\,\mathrm{Hz}$.
Its peak strain amplitude nevertheless decreased from approximately
$9.46\times10^{-21}$ in the earliest block to $4.28\times10^{-21}$
on-source, corresponding to about an $80\%$ decline in line power.
A log-linear amplitude trend fitted exclusively to the four training blocks
had an amplitude e-folding time of approximately $1389\,\mathrm{s}$.
It predicted the later fitted line powers within about $6\%$, and the
on-source fitted line power within about $9\%$. This demonstrates amplitude
evolution in the cached strain; it does not identify the physical cause or
establish a transient-glitch classification.

The forward prediction used a sinusoid with that training-only frequency
and amplitude trend, including its within-block envelope, plus an estimated
local continuum. Its exact phase-averaged response was computed through
the symmetric Tukey window used by the on-source FFT. This did not convolve
an already-windowed PSD a second time. Synthetic-tone recovery,
window-weighted power conservation, and agreement with the pipeline's
on-source FFT periodogram were checked.

For the actual on-source Tukey window, the observed-to-predicted integrated
power ratios were $0.913$ in the core and $1.044$ in the skirts. The core
was defined within $0.25\,\mathrm{Hz}$ of the training frequency; the skirts
were $0.25$--$4\,\mathrm{Hz}$ away. Across the four held-out off-source
blocks these ratios ranged from $0.938$ to $1.004$ in the core and from
$1.022$ to $1.142$ in the skirts. In contrast, the stationary-amplitude
prediction gave an on-source core ratio of only $0.269$: matching the window
alone did not account for the change in line strength.

Within the local cache, this establishes a finite-window response and
nonstationary line amplitude. It does not by itself establish their share
of the cross-release Welch--BayesLine discrepancy. Agreement in
integrated power is not a complete calibration: the on-source mean
per-bin skirt-power ratio was still about $1.69$, versus a median of
$0.95$, indicating remaining localized excess. A Hann-window diagnostic
exposed additional narrow structures, leaving skirt power about $14$ times
the single-line-plus-continuum prediction. Parameter uncertainty,
additional lines, and frequency-bin covariance remain unmodelled. Local
core-conditioned fits are explanatory controls, and the reported diagonal
noise-score shifts are not model-selection Bayes factors. The BayesLine PSD
was displayed only as an on-source reference, not as an independently
trained held-out model. Production windows, PSD settings, and queued jobs
were unchanged.

\subsection{Matched No-Glitch BayesLine Control}

A diagnostic on 30 October 2026 matched the nominal Welch settings to the
archived \texttt{noiseOnly}, \texttt{noClean}, \texttt{noWaveletPrior}
BayesLine run. Its analysis interval was GPS
$[1187008610.4,1187008738.4]$, with $4096\,\mathrm{Hz}$ sampling,
a $128\,\mathrm{s}$ FFT, $\Delta f=0.0078125\,\mathrm{Hz}$, and the
$23$--$2048\,\mathrm{Hz}$ band. A $0.4\,\mathrm{s}$ edge taper,
corresponding to Tukey $\alpha=0.00625$, reproduces the logged window
correction to its printed precision. The separate configuration option
\texttt{window=1.0} was not interpreted as a Tukey alpha.

One complete segment produces a single periodogram, not a variance-reduced
Welch average. Its normalization was checked against a direct FFT and
window-weighted Parseval identity. Raw posterior PSDs were converted by
$2/128$ and checked against the archived physical-units export; the same
conversion was applied to the archived starting PSD. BayesWavePost omits
Nyquist and rounds frequency labels, so ordered rows were checked against
the exact starting-PSD grid instead of interpolating the rounded labels.

The nominally matched local strain still disagreed strongly with the
BayesLine PSD. However, the archived no-glitch Fourier residual export
revealed a more important distinction: it was not numerically equivalent
to the local strain. Integrated data-to-BayesLine power ratios were
\[
\begin{array}{lrr}
 & \text{Local cache} & \text{Archived Fourier export}\\
\text{Core: within }0.25\,\mathrm{Hz} & 1312.61 & 0.693\\
\text{Skirts: }0.25\text{--}4\,\mathrm{Hz} & 10.993 & 0.900\\
\text{Background: }5\text{--}8\,\mathrm{Hz} & 1.036 & 0.971
\end{array}
\]
around the fitted $1875.3920\,\mathrm{Hz}$ line. The archived export had
about $1900$ times less integrated core power than the local cache.
Symmetric versus periodic Tukey windows changed local core power by only
about $3.1\times10^{-7}$ fractionally. The nearest local sample starts
$0.4004$ samples earlier than the requested GPS time; adjacent-sample
sensitivity controls were also recorded.

The executed run flags disable glitch, signal, and CBC components, making
its Fourier residual a useful processed-data control. Nevertheless, it is
an exported residual product, not an independently verified copy of the
unconditioned frame strain. These results therefore locate an unresolved
data-provenance, conditioning, or export-consistency issue; they do not
identify its cause, demonstrate deliberate line removal, or establish that
BayesLine intrinsically misses the line. The archive-only core and skirt
ratios are of order one, but this is an in-sample diagnostic, not noise-model
calibration. The earlier conclusion that matching nominal settings alone
tests estimator differences must be qualified by this processed-data check.

The full $1024\,\mathrm{s}$ initialization context was not reproduced:
its first $2.075\,\mathrm{s}$ are absent from the local cache. No new
BayesWave or CBC inference jobs were launched, and production settings were
unchanged. Numerical checks, source hashes, arrays, and plots are saved by
the \texttt{--compare-no-glitch} diagnostic in its separate
\texttt{matched\_no\_glitch} output directory.

\subsection{Original-Input Release Identification and L1 Glitch Status}

On 1 November 2026 the actual frame input was checked rather than relying
on the post-processing export or nominal estimator settings. The archived
H1/L1/V1 frame-cache records name the $4096\,\mathrm{Hz}$ GWTC-1
GW170817 event-v3 files beginning at GPS $1187006835$, with duration
$4096\,\mathrm{s}$. SHA256 hashes of all three local frames matched their
current official GWOSC download bytes exactly. Duplicate local frame
copies also matched. This authenticates the available original input;
it is not a recovery of the historical download log.

The local sampler cache was separately compared with the official
O1\_O2-Preliminary GW170817 event-v2 cleaned HDF5 products beginning at
GPS $1187007040$. \emph{Every one of the $4{,}784{,}128$ aligned samples
per detector matched exactly}, over the complete cached
$1168\,\mathrm{s}$ interval
$[1187007716.4746094,1187008884.4746094)$. All H1/L1/V1 cache copies
checked against the local historical archive also had identical file
hashes. Thus the cached strain is identified as the published v2 release,
not as an unspecified O2 bulk product. The collection code's
\texttt{fetch\_open\_data(version=2)} is consistent with this finding,
but the identification rests on sample equality, not the version argument.

For the nominal $128\,\mathrm{s}$ no-glitch control, the first frame
sample at or after the requested start is GPS $1187008610.4001465$.
Using the same aligned local samples, symmetric Tukey taper and
periodogram normalization gives original-frame/local-cache integrated
V1 power ratios
\[
\begin{array}{lr}
\text{Core: within }0.25\,\mathrm{Hz}\text{ of }1875.3920\,\mathrm{Hz}
 & 0.00053034\\
\text{Skirts: }0.25\text{--}4\,\mathrm{Hz} & 0.0822164\\
\text{Broadband: }100\text{--}1500\,\mathrm{Hz} & 0.936858.
\end{array}
\]
The original frame already has approximately $1886$ times less core power.
The archived V1 Fourier export closely reproduces that frame: the relative
complex residual is approximately $0.00119$, with fitted complex scale
$0.99999908+0.00000081i$. Its integrated power divided by the physical
frame periodogram is approximately $0.996093$, consistent with the
window-energy factor $\langle w^2\rangle$ omitted from the export's
convention. A one-sample difference between nearest-sample and
first-at-or-after-start extraction materially affects the complex phases,
but not the orders-of-magnitude power distinction.

\textbf{The large V1 discrepancy is therefore present in the input strain
before BayesLine fitting.} It is not evidence that BayesLine intrinsically
misses the strong line seen in a different release. The successful complex
match also rules out a wholly independent simulated-noise realization as
an adequate explanation for this export, despite the simulated-channel
labels in the BayesWavePost log. The exact remaining conditioning,
calibration and downsampling differences have not been isolated, and the
precise historical post-processing overwrite sequence is not inferred.

The known strong GW170817 glitch is a separate L1 issue. GWOSC documents
its time as GPS $1187008881.389$, about $1.1\,\mathrm{s}$ before merger,
and a brief digital-to-analog converter saturation of less than
$5\,\mathrm{ms}$. GWTC-1 event-v3 strain retains the glitch; the earlier
v2 release is explicitly documented as glitch-subtracted. Since the full
cached L1 samples exactly match v2, \emph{the local sampler cache already
uses the published glitch-subtracted L1 data}. A common diagnostic
$30$--$1000\,\mathrm{Hz}$ zero-phase bandpass shows the transient in the
original v3 frame and its suppression in the cache. This plot is a display
control, not a reconstruction-validation test or a new likelihood.
Unfiltered raw maxima alone are misleading because the v2 H1/L1 products
have documented excess noise below $20\,\mathrm{Hz}$. The v2 H1
last-$150\,\mathrm{s}$ cleaning-window warning concerns approximately
GPS $1187008938$ onward, outside the cached interval checked here.

The off-source BayesWave control ends at GPS $1187008738.4$, before the
L1 glitch. Off-source PSD cleaning neither explains the cross-detector
V1 difference nor establishes on-source glitch subtraction. In this case,
the cache's subtraction status is established independently by its release
content. Special glitch-subtracted products described by GWOSC also differ
in calibration and glitch model between publications; they must not be
interchanged merely because all are called cleaned data.

There is an additional frequency-band caveat. GWOSC's O2 technical notes
recommend the $16384\,\mathrm{Hz}$ products for studies around
$1700\,\mathrm{Hz}$ and above because the $4096\,\mathrm{Hz}$
anti-alias response corrupts the Nyquist region. The $1875\,\mathrm{Hz}$
diagnostic cannot be promoted to physical high-frequency calibration
without explicit release and filter checks. No frequency cut, data
replacement, gating, PSD regeneration or new inference job was applied.
The frozen calibration bundle remains a legacy mixed-release fixed-PSD
comparison; its historical unresolved-glitch caveat is now qualified by
this independent sample audit, not rewritten in its frozen configuration.

Earlier tables and score snapshots are preserved. The historical
``at least $88\%$'' localization claim above was not reproduced by a
later direct net-score integration, which gave approximately $67\%$
for the $1871$--$1879\,\mathrm{Hz}$ share of the net V1 shift; these
definitions or snapshots must not be combined. Neither those score tests
nor the present release-identification checks validate a detector-consistent
PSD likelihood or a completed GW170817 recovery. A fair estimator
comparison requires PSD estimation and on-source strain from a matched
release with explicit conditioning and a defensible analysis band.

The independent original-input audit saves official API snapshots,
file and sample hashes, aligned spectra, ratios and raw/bandpassed glitch
plots in a fresh timestamped directory. All $16$ hashed source/config/input
files in the frozen calibration bundle were checked without modification.
Sources are the GWOSC GWTC-1 documentation, the GW170817 event-v2/v3
release records, the dedicated GW170817 strain documentation and the O2
technical details; URLs and downloaded-file hashes are recorded in the
audit report. No legacy inference results were replaced.

\subsection{Top-Candidate Cross-Checks}

The same detector- and frequency-resolved normalized-score diagnostic was
applied to the two leading subthreshold candidates, GRB190728271 and
GRB240715239. Unlike GW170817, their discrepancies do not support a common
V1 feature near $1875\,\mathrm{Hz}$. Table~\ref{tab:top-candidate-psd-scores}
uses the sign convention $\ell^{\rm BL}-\ell^{\rm Welch}$; a positive value
therefore indicates a higher normalized fixed-PSD score for BayesLine.

\begin{table}[ht]
\centering
\small
\caption{Detector-resolved normalized fixed-PSD score differences for the
top candidates. The CBC entries use the conditioned maximum-log-likelihood
posterior waveform with the common coalescence phase profiled separately for
each PSD. These are predictive-score diagnostics, not Bayes factors.}
\label{tab:top-candidate-psd-scores}
\begin{tabular}{llrr}
\toprule
\textbf{Event} & \textbf{Detector} & \textbf{Noise-only}
& \textbf{CBC residual}\\
& & \textbf{$\ell^{\rm BL}-\ell^{\rm Welch}$} &
\textbf{$\ell^{\rm BL}-\ell^{\rm Welch}$}\\
\midrule
GRB190728271 & H1 & $+3031.38$ & $+3032.01$\\
GRB190728271 & L1 & $+4583.71$ & $+4581.30$\\
GRB190728271 & Network & $+7615.09$ & $+7613.31$\\
\addlinespace
GRB240715239 & L1 & $+4379.88$ & $+4377.77$\\
GRB240715239 & V1 & $+7616.29$ & $+7622.69$\\
GRB240715239 & Network & $+11996.16$ & $+12000.46$\\
\bottomrule
\end{tabular}
\end{table}

GRB190728271 was observed only in H1 and L1, so a V1 masking test is not
defined for that event. Its score difference is distributed across both
detectors. The largest $1\,\mathrm{Hz}$ contributions occur around
$1461$--$1465\,\mathrm{Hz}$ in H1 and $1022$--$1027\,\mathrm{Hz}$ in L1,
showing that its PSD sensitivity cannot be attributed to the GW170817 V1
feature.

For GRB240715239, dropping V1 leaves the L1-only normalized difference of
$+4379.88$ nats, rather than reversing the sign as it did for GW170817. V1 is
important, contributing $+7616.29$ nats, but L1 also contributes materially.
The targeted V1 interval $1871$--$1879\,\mathrm{Hz}$ contributes only
$+232.12$ nats, or $3.05\%$ of the V1 noise-only difference. Its largest V1
score differences instead lie predominantly between $31$ and
$64\,\mathrm{Hz}$. Thus, although both candidates are sensitive to the fixed
PSD choice, neither has evidence for the specific GW170817-like
$1875\,\mathrm{Hz}$ V1 discrepancy.

As above, detector removal or frequency masking changes the likelihood.
These results localize the fixed-PSD sensitivity, but they do not yield
masked CBC evidences or conditioned-versus-unconditioned Bayes factors.
Those quantities require detector- and mask-consistent likelihood evaluations
and nested-sampling reruns.

The conclusion is therefore methodological. Fixed-PSD evidence differences
should be treated as PSD-sensitivity diagnostics. They motivate inspections
of the affected detector and frequency range, repeat runs using posterior PSD
draws, off-source control studies, and software injections. A definitive
comparison would propagate PSD uncertainty and model non-stationary lines,
window-induced frequency correlations, and transient glitches jointly with the
CBC signal. Finally, removing V1 after the fact permits only an H1/L1
noise-score diagnostic; it does not produce an H1/L1 CBC Bayes factor or
evidence without rerunning the likelihood and nested sampler for that detector
network.

\subsection{Detector- and Frequency-Resolved Likelihood Digest}

For a detector $I$ and frequency bin $k$, define the normalized fixed-PSD
Whittle contribution for a waveform $h$ as
\[
\ell_{I,k}(h;S_n) =
\log\!\left(\frac{2\Delta f}{\pi S_{n,I}(f_k)}\right)
- \frac{2\Delta f\lvert\tilde d_{I,k}-\tilde h_{I,k}\rvert^2}
{S_{n,I}(f_k)}.
\]
The detector score and network score are respectively
\[
\ell_I(h;S_n) = \sum_{k\in[23,2048]~\mathrm{Hz}}\ell_{I,k}(h;S_n),
\qquad
\ell_{\rm net}(h;S_n) = \sum_I\ell_I(h;S_n).
\]
For the noise-only diagnostic, set $h=0$. For the CBC diagnostic, $h$ is a
specified posterior waveform, with a common coalescence phase profiled
separately for each fixed PSD. These are predictive scores for a selected
waveform; they are not nested-sampling evidences and should not replace the
evidence table above.

For direct comparison of the two PSD curves, the detector- and bin-resolved
differences are
\[
\delta\ell_I(h) = \ell_I(h;S_n^{\rm BL})-
\ell_I(h;S_n^{\rm Welch}),
\qquad
\delta\ell_{I,k}(h) =
\ell_{I,k}(h;S_n^{\rm BL})-
\ell_{I,k}(h;S_n^{\rm Welch}),
\]
with $\delta\ell_I=\sum_k\delta\ell_{I,k}$. A positive value means that
BayesLine assigns a higher normalized score; a negative value means that
Welch assigns a higher normalized score. This convention makes the V1
discrepancy transparent: the narrow bin near $1875\,\mathrm{Hz}$ gives a
large negative $\delta\ell_{\mathrm{V1},k}$ when the BayesLine PSD is low
relative to the observed on-source line power.

\begin{table}[ht]
\centering
\small
\caption{Detector-resolved format for the normalized Whittle audit. Each row
is obtained by summing the bin contributions in the stated detector over the
analysis band. The network row is the sum of detector rows. Report the values
from the audit output for the selected event and posterior waveform.}
\begin{tabular}{llrrrr}
\toprule
\textbf{Detector} & \textbf{Waveform} & \textbf{$\ell_I^{\rm Welch}$}
& \textbf{$\ell_I^{\rm BL}$} & \textbf{$\delta\ell_I$}
& \textbf{Interpretation}\\
\midrule
H1 & $h=0$ & \textit{audit value} & \textit{audit value}
& \textit{BL $-$ Welch} & noise-only score\\
L1 & $h=0$ & \textit{audit value} & \textit{audit value}
& \textit{BL $-$ Welch} & noise-only score\\
V1 & $h=0$ & \textit{audit value} & \textit{audit value}
& \textit{BL $-$ Welch} & inspect narrow lines\\
Network & $h=0$ & \textit{sum} & \textit{sum}
& \textit{sum} & detector sum\\
\addlinespace
H1 & $h=h_{\rm prof}$ & \textit{audit value} & \textit{audit value}
& \textit{BL $-$ Welch} & CBC residual score\\
L1 & $h=h_{\rm prof}$ & \textit{audit value} & \textit{audit value}
& \textit{BL $-$ Welch} & CBC residual score\\
V1 & $h=h_{\rm prof}$ & \textit{audit value} & \textit{audit value}
& \textit{BL $-$ Welch} & CBC residual score\\
Network & $h=h_{\rm prof}$ & \textit{sum} & \textit{sum}
& \textit{sum} & detector sum\\
\bottomrule
\end{tabular}
\end{table}

Rather than tabulating every $1/128\,\mathrm{Hz}$ Fourier bin, which would
obscure the result, the per-bin audit lists the bins or narrow bands with the
largest absolute score differences. The confirmed V1 feature can be reported
compactly as follows:
\begin{table}[ht]
\centering
\small
\caption{Largest frequency-resolved contributions to the fixed-PSD score
difference. Values are computed from $\delta\ell_{I,k}$; neighbouring bins
may be grouped only when their summed contribution is reported.}
\begin{tabular}{llrrrl}
\toprule
\textbf{Event} & \textbf{Detector} & \textbf{Frequency/band [Hz]}
& \textbf{$\delta\ell_{I,k}(0)$} & \textbf{$\delta\ell_{I,k}(h_{\rm prof})$}
& \textbf{Feature}\\
\midrule
GRB170817529 & V1 & $1871$--$1879$ &
$\sum_{k\in\mathrm{band}}\delta\ell_{\rm V1,k}(0)
\leq-0.88\lvert\delta\ell_{\rm V1}(0)\rvert$
& \textit{audit value} & narrow feature and leakage skirts\\
\textit{event} & \textit{IFO} & \textit{largest remaining bin/band}
& \textit{audit value} & \textit{audit value} & \textit{diagnostic label}\\
\textit{event} & \textit{IFO} & \textit{largest remaining bin/band}
& \textit{audit value} & \textit{audit value} & \textit{diagnostic label}\\
\bottomrule
\end{tabular}
\end{table}

The two tables should always be read together. A large detector total can be
broad-band, or it can be dominated by one narrow instrumental feature. The
per-bin table distinguishes these cases and prevents a narrow V1 line from
being misdescribed as broad astrophysical residual power.